\documentclass[11pt,oneside,openany]{book}
\usepackage[T1]{fontenc}
\usepackage[utf8]{inputenc}
\usepackage{amsmath,amsthm,mathtools}
\usepackage{mathpazo,amssymb}
\usepackage[paperwidth=7.5in,paperheight=10in,inner=0.80in,outer=0.80in,top=0.82in,bottom=0.78in,headheight=14pt,headsep=0.22in,footskip=0.35in]{geometry}
\usepackage{microtype,graphicx,booktabs,array,longtable,tabularx}
\usepackage[table]{xcolor}
\usepackage{tikz}
\usetikzlibrary{arrows.meta,positioning,calc}
\usepackage{titlesec,fancyhdr,enumitem,caption,float,needspace,etoolbox,listings}
\usepackage[most]{tcolorbox}
\usepackage{hyperref,bookmark}
\definecolor{Ink}{HTML}{172A3A}
\definecolor{Accent}{HTML}{286E75}
\definecolor{Soft}{HTML}{F0F5F5}
\definecolor{Quiet}{HTML}{62727C}
\definecolor{CoverLight}{HTML}{CDDEE2}
\hypersetup{colorlinks=true,linkcolor=Accent,citecolor=Accent,urlcolor=Accent,pdftitle={Realization Boundaries: Finite Grammar, Infinite Groups, and the Limits of Reconstruction},pdfauthor={Komorebi P.},pdfsubject={Realization Calculus, abelian groups, profinite completion, computable structures},bookmarksnumbered=true}
\urlstyle{same}
\renewcommand{\bibname}{References}
\setlength{\parindent}{1.05em}
\setlength{\parskip}{3.5pt plus 0.8pt}
\linespread{1.045}
\setlength{\emergencystretch}{2em}
\raggedbottom
\setcounter{tocdepth}{1}
\setcounter{secnumdepth}{2}
\titleformat{\chapter}[display]{\normalfont\color{Ink}}{\sffamily\small\color{Accent}\MakeUppercase{\chaptertitlename}\ \thechapter}{9pt}{\fontsize{25}{29}\selectfont\raggedright\hyphenpenalty=10000\exhyphenpenalty=10000}
\titlespacing*{\chapter}{0pt}{4pt}{20pt}
\titleformat{\section}{\large\bfseries\color{Ink}}{\thesection}{0.7em}{}
\titleformat{\subsection}{\normalsize\bfseries\color{Ink}}{\thesubsection}{0.7em}{}
\titlespacing*{\section}{0pt}{16pt}{7pt}
\titlespacing*{\subsection}{0pt}{11pt}{5pt}
\pagestyle{fancy}
\fancyhf{}
\fancyhead[L]{\sffamily\fontsize{8}{10}\selectfont\color{Quiet}\nouppercase{\leftmark}}
\fancyhead[R]{\sffamily\fontsize{8}{10}\selectfont\color{Quiet}REALIZATION CALCULUS}
\fancyfoot[C]{\small\thepage}
\renewcommand{\headrulewidth}{0.3pt}
\renewcommand{\chaptermark}[1]{\markboth{\thechapter\quad #1}{}}
\fancypagestyle{plain}{\fancyhf{}\fancyfoot[C]{\small\thepage}\renewcommand{\headrulewidth}{0pt}}
\captionsetup{font=small,labelfont={bf,color=Accent},justification=raggedright,singlelinecheck=false,skip=8pt}
\setlist{itemsep=3pt,topsep=6pt,parsep=0pt,leftmargin=1.6em}
\newtheoremstyle{research}{9pt}{9pt}{\itshape}{} {\bfseries\color{Ink}}{.}{0.55em}{}
\theoremstyle{research}
\newtheorem{theorem}{Theorem}[chapter]
\newtheorem{proposition}[theorem]{Proposition}
\newtheorem{lemma}[theorem]{Lemma}
\newtheorem{corollary}[theorem]{Corollary}
\theoremstyle{definition}
\newtheorem{definition}[theorem]{Definition}
\newtheorem{example}[theorem]{Example}
\newtheorem{remark}[theorem]{Remark}
\newtheorem{boundary}[theorem]{Boundary}
\newtheorem{question}[theorem]{Question}
\newtcolorbox{keybox}[1][]{breakable,colback=Soft,colframe=Accent,boxrule=0.5pt,arc=0pt,left=11pt,right=11pt,top=10pt,bottom=10pt,before skip=12pt,after skip=12pt,#1}
\newcommand{\N}{\mathbb N}
\newcommand{\Z}{\mathbb Z}
\newcommand{\Q}{\mathbb Q}
\newcommand{\F}{\mathbb F}
\newcommand{\R}{\mathbb R}
\newcommand{\Tor}{\operatorname{Tor}}
\newcommand{\Hom}{\operatorname{Hom}}
\newcommand{\End}{\operatorname{End}}
\newcommand{\Aut}{\operatorname{Aut}}
\newcommand{\supp}{\operatorname{supp}}
\newcommand{\rad}{\operatorname{rad}}
\newcommand{\den}{\operatorname{den}}
\newcommand{\Hull}{\operatorname{Hull}}
\newcommand{\Sub}{\operatorname{Sub}}
\newcommand{\srank}{\operatorname{srank}}
\newcommand{\trank}{\operatorname{trank}}
\newcommand{\fin}{\operatorname{fin}}
\newcommand{\Obs}{\mathcal F}
\newcommand{\one}{\mathbf 1}
\newcommand{\eqfin}{=^{*}}
\newcommand{\neqfin}{\ne^{*}}
\newcommand{\primes}{\mathcal P}
\newcommand{\bits}{\{0,1\}^{\N}}
\newcommand{\proofend}{\hfill$\square$}
\newcommand{\chapterintro}[1]{\noindent\textit{#1}\par\vspace{5pt}}
\newcommand{\figfile}[1]{../figures/#1}
\lstset{basicstyle=\ttfamily\footnotesize,columns=fullflexible,keepspaces=true,breaklines=true,frame=single,rulecolor=\color{Quiet!35},backgroundcolor=\color{Soft},xleftmargin=4pt,xrightmargin=4pt,aboveskip=10pt,belowskip=10pt,showstringspaces=false,upquote=true}
\begin{document}
\begin{titlepage}
\thispagestyle{empty}
\begin{tikzpicture}[remember picture,overlay]
\fill[Ink] (current page.south west) rectangle (current page.north east);
\draw[Accent,line width=2pt] ([xshift=0.82in,yshift=-1.28in]current page.north west)--([xshift=-0.82in,yshift=-1.28in]current page.north east);
\node[anchor=north west,text=CoverLight,font=\sffamily\small,align=left] at ([xshift=0.82in,yshift=-0.85in]current page.north west) {REALIZATION CALCULUS};
\node[anchor=north west,text=white,text width=5.7in,align=left,font=\fontsize{40}{45}\selectfont] at ([xshift=0.82in,yshift=-1.90in]current page.north west) {Realization\\Boundaries};
\node[anchor=north west,text=CoverLight,text width=5.6in,align=left,font=\fontsize{17}{23}\selectfont] at ([xshift=0.82in,yshift=-3.65in]current page.north west) {Finite Grammar, Infinite Groups,\\and the Limits of Reconstruction};
\node[anchor=west,text=white,font=\fontsize{22}{28}\selectfont] at ([xshift=0.82in,yshift=2.65in]current page.south west) {$D\ \subseteq\ G_A\ \subseteq\ P$};
\node[anchor=north west,text=CoverLight,text width=5.4in,align=left,font=\fontsize{11}{16}\selectfont] at ([xshift=0.82in,yshift=2.20in]current page.south west) {One ambient completion.\\Different algebraic realizations.\\Precisely measured information.};
\node[anchor=south west,text=white,font=\large] at ([xshift=0.82in,yshift=0.82in]current page.south west) {Komorebi P.};
\end{tikzpicture}
\end{titlepage}
\frontmatter
\chapter*{Preface}
\addcontentsline{toc}{chapter}{Preface}
\markboth{Preface}{}
An infinite object can be described by finite rules without being generated by finitely many elements. It can be defined by a short mathematical sentence without admitting computable operations. It can be separated point by point by finite observations without being recoverable from the entire unmarked system of those observations. These distinctions are the subject of this book.

Realization Calculus studies the passage from specified data and admissible operations to an actual mathematical object. Its boundary questions are therefore not external qualifications. They ask which information a realization requires, which operations preserve that information, and which observations discard it. Here these questions are examined in a concrete universe of abelian groups. The constructions are elementary enough to calculate with, yet expressive enough to carry every Turing degree.

The starting object is the direct sum of one cyclic group of each prime order. Its finite stages are cyclic. Its completion is topologically cyclic. The direct sum itself nevertheless needs infinitely many seeds even when every group endomorphism is allowed. Its subgroup lattice then supplies a transparent first model of finite description and effective presentation.

The central construction is more stringent. For every set $A\subseteq\N$, a countable group $G_A$ is placed between that direct sum $D$ and its product completion $P$. All groups have the same torsion subgroup, the same rational quotient, and the same full unmarked finite-quotient system. Nevertheless, their abstract isomorphism types are classified by $A$ modulo finite changes, and an oracle computes a copy precisely when it computes $A$.

\begin{keybox}
\textbf{The central separation.} The compatible finite quotients determine the ambient completion. They do not determine which dense subgroup is the intended original realization. In this family the omitted choice has an exact divisibility signature and an exact computational cost.
\end{keybox}

\section*{How to read the book}
Chapters 1--3 separate the realization contracts and develop the first obstruction examples. Chapters 4--8 give the main construction and its proofs. Chapter 4 also describes the general ambient mechanism: rational subspaces of $P/D$ correspond to pure intermediate groups. Chapters 9--10 extract the reconstruction and grammar consequences. Chapter 11 provides calculations that can be followed by hand and reproduced in Python. Chapter 12 returns to the interpretation in Realization Calculus. The appendices contain implementation guidance, a theorem map, a formal proof guide, and a glossary.

Only elementary group theory is needed at the start. The later arguments introduce profinite completion and the relevant computability notions before using them. Readers interested chiefly in the common-completion theorem can begin with Chapters 4--6. Readers interested in its effective meaning should then read Chapters 7--8, since the distinction between a marked and an unmarked group is essential to the decoding proof.

\section*{Evidence and mathematical context}
The main results have ordinary proofs in the text and a Lean companion covering the numbered contracts. The effective statements use explicit oracle and presentation semantics. Separate Python programs use exact rational and modular arithmetic to test finite instances and reject selected false formulas. These finite tests support the implementation. The infinite conclusions are established by the mathematical proofs and the formal companion.

The prime-spectrum coding in Chapter 3 has an explicit classical antecedent in Montalb\'an's text \cite{Montalban}. Degrees of isomorphism types, pure subgroups, cocycle extensions, and profinite completion belong to established theories \cite{Calvert,Fuchs,Lupini,Wilkes}. The common-completion family is proved directly rather than attributed to an unverified source. Historical priority for that combined construction has not been established. The contribution presented here is a self-contained theorem development organized around precise realization contracts, not a claim that the classical ingredients originate in this work.

\section*{Conventions that prevent false conclusions}
An unmarked finite-quotient system does not contain the original source group or a distinguished coherent source element. An infinite parameter is not made finite by assigning it a single letter. A computable copy is not the same as an enumerable list of generators and relators. A nonzero divisible quotient does not imply a divisible subgroup. These distinctions appear where they are used, rather than being left to a general disclaimer.

\begingroup\setlength{\parskip}{0pt}\tableofcontents\endgroup
\chapter*{Notation and principal results}
\addcontentsline{toc}{chapter}{Notation and principal results}
\markboth{Notation and principal results}{}
All groups are written additively unless a finite target group is explicitly allowed to be nonabelian. The natural numbers start at zero. Primes are indexed by
\[
p_0=2,\quad p_1=3,\quad p_2=5,\quad p_3=7,\quad\ldots.
\]
\begin{longtable}{@{}p{0.23\textwidth}p{0.71\textwidth}@{}}
\toprule
Symbol & Meaning\\\midrule
$C_p=\F_p$ & The additive cyclic group of prime order $p$.\\
$D=\bigoplus_p C_p$ & Finite-support prime-coordinate vectors.\\
$P=\prod_p C_p$ & All prime-coordinate vectors, with product topology.\\
$q:P\to V$ & The quotient map to $V=P/D$.\\
$D_B$ & The subgroup supported on prime indices in $B\subseteq\N$.\\
$H_W=q^{-1}(W)$ & The intermediate group associated to a rational subspace $W\le V$.\\
$S_A$ & The prime set $\{p_{2n+\one_A(n)}:n\in\N\}$.\\
$u_A$ & The coordinate indicator of $S_A$ in $P$.\\
$G_A$ & The group $q^{-1}(\Q q(u_A))$.\\
$s_A(r)$ & The specified rational lift of $r q(u_A)$ to $P$.\\
$c_A(r,t)$ & The finite cocycle $s_A(r)+s_A(t)-s_A(r+t)$.\\
$e_p$ & The vector supported at $p$ with value one.\\
$p\mid_G g$ & There exists $y\in G$ with $py=g$.\\
$A\eqfin B$ & The symmetric difference $A\triangle B$ is finite.\\
$\rad(m)$ & The product of the distinct prime divisors of $m$, with $\rad(1)=1$.\\
$\widehat G$ & The profinite completion, not the Pontryagin dual.\\
$A\le_T X$ & Membership in $A$ is computable using oracle $X$.\\
$K_{\mathfrak G}(G)$ & Least admissible finite description length, or $+\infty$.\\
\bottomrule
\end{longtable}

\noindent The main statements have the following form:
\begin{align*}
\Tor(G_A)&=D,&G_A/D&\cong\Q,\\
G_A/mG_A&\cong C_{\rad(m)},&\widehat{G_A}&\cong P,\\
G_A\cong G_B&\iff A\eqfin B,&&\\
G_A\text{ has an }X\text{-computable copy}&\iff A\le_T X.&&
\end{align*}
The extension by $D$ does not split. Distinct eventual selector classes do not embed in one another. All of these conclusions are proved within the book, with the exact hypotheses retained.

\mainmatter
\chapter{What a realization must supply}
\label{ch:contracts}
\chapterintro{Before constructing a boundary, one must specify what is allowed to cross it. A finite seed family, a finite description, and an effective presentation impose different requirements.}

\section{Four questions, not one}
The expression ``finitely realizable'' is incomplete unless it names a contract. A group may be infinite but cyclic. A group may require infinitely many algebraic generators but have a short recursive construction. A description may select an exact group through a noncomputable set. Finally, a topology may allow limits that are not present in the underlying algebraic construction.

The integer group $\Z$ illustrates only the first distinction. One generator and its integer multiples produce infinitely many elements. The groups developed below separate the other distinctions as well. Their purpose is not to make infinite objects mysterious. It is to account for the information that a construction actually consumes.

\begin{definition}[Operator-seed realization]
Let $M$ be an abelian group and let $\mathcal O\subseteq\End(M)$ be a declared collection of endomorphisms. For $S\subseteq M$, its operator hull is the subgroup generated by all $T(s)$, where $s\in S$ and $T$ is a finite composition of members of $\mathcal O$, including the identity. A finite operator-seed realization requires a finite $S$ whose hull is $M$.
\end{definition}
The allowed addition and inverse operations are part of the group structure. An arbitrary target-dependent map is not automatically an admissible endomorphism. In particular, one cannot evade a generation obstruction by introducing a map that simply returns the desired answer.

If all endomorphisms are allowed, this is module generation over the unital ring $\End(M)$. We write $\srank_{\End(M)}(M)$ for its least generating cardinal. An infinite value can therefore express a strong obstruction even when the operator collection itself is unrestricted.

\begin{definition}[Finite semantic description]
A finite-description grammar consists of a countable alphabet, a class of valid finite strings, and a semantics assigning each valid string one group isomorphism type. The semantics need not be computable. Any permitted parameters must have a declared coding convention.
\end{definition}
This definition allows a finite sentence to specify a noncomputable group. It does not count an arbitrary real or subset of $\N$ as a free symbol. If such a parameter is admitted, its information must be recorded separately.

\begin{definition}[Computable copy]
A computable copy of a countable group is an isomorphic group on a computably named domain, with decidable equality and computable group operations. An $X$-computable copy permits an oracle $X$ in those computations. Infinite groups may equivalently be presented on $\N$.
\end{definition}
A presentation by generators and relators is a different object. Enumerating the relators need not decide whether two words represent the same element. The effective presentations in this book always include decidable equality in the named copy.

\begin{definition}[Topological generation]
For a topological group $M$, a subset $S$ topologically generates $M$ when the closure of the subgroup generated by $S$ is all of $M$. Its least possible size is denoted $\trank(M)$. An operator version takes the closure of the operator hull instead.
\end{definition}
The closure is taken in the stated topology. A change of topology is a change of contract, not an innocuous notational adjustment.

\section{Finite descriptions and infinite information}
There are countably many finite strings over a countable alphabet. Consequently a grammar with single-valued semantics realizes at most countably many group isomorphism types. This observation requires neither an algorithm for checking validity nor an algorithm for evaluating the semantics.

\begin{proposition}[Countable-code bound]
\label{thm:codecount}
Let $\mathfrak G$ be a finite-description grammar as above. Then its class of realized group isomorphism types has cardinal at most $\aleph_0$.
\end{proposition}
\begin{proof}
For each length $n$, the set of strings of length $n$ is countable. The union over $n\in\N$ is countable. The valid strings form a subset, and a single-valued image of a countable set is countable.
\end{proof}
This is a boundary of description, not of existence. It also says nothing about whether a particular named group lies outside a chosen grammar. The later constructions supply structural examples and an explicit relative diagonal argument.

For quantitative description length, fix an injective encoding of descriptions by finite binary strings. A symbol from a countable alphabet must carry its finite code rather than count as a free unit of information. This encoding convention does not require the semantic evaluation to be computable.

\begin{definition}[Grammar-relative description complexity]
For a fixed binary coding of $\mathfrak G$, let $\lvert d\rvert$ be its bit length and set
\[
K_{\mathfrak G}(G)=\inf\{\lvert d\rvert:\ d\text{ is valid and }\operatorname{Realize}_{\mathfrak G}(d)\cong G\},
\]
with value $+\infty$ when there is no such description.
\end{definition}
The value depends on the grammar and the coding. A fixed finite description can be long without being impossible. Conversely, a group with $K_{\mathfrak G}(G)=+\infty$ may have a short description in an expanded language.

\begin{boundary}[No uncharged infinite parameters]
The expression $G=q^{-1}(\Q q(u))$ is short. If $u$ is an arbitrary infinite coordinate vector, that expression is not a finite specification of $G$ without a specification of $u$. Counting its surrounding symbols does not count the information in the parameter.
\end{boundary}
This boundary will become exact: the computational information in a selected $u_A$ is precisely what is needed to compute the corresponding group.

\section{A reusable support obstruction}
A finite-seed failure can sometimes be certified without examining every possible operator word. The packet structure supplies a witness.

\begin{theorem}[Reachability obstruction]
\label{thm:reach}
Let $M=\bigoplus_{i\in I}M_i$ be a direct sum of abelian groups with coordinate projections $P_i$. Suppose every primitive operator $T$ is an endomorphism and its transfer block $P_jTP_i$ vanishes unless $i\to j$ is an allowed edge of a directed graph $E$. If $F$ is the union of the supports of a seed family $S$, then
\[
\Hull_{\mathcal O}(S)\subseteq
\bigoplus_{j\in\operatorname{Reach}_E(F)}M_j.
\]
If $M_j\ne0$ and $j$ is not reachable, $P_j$ annihilates the entire hull but not $M$.
\end{theorem}
\begin{proof}
An element of the direct sum has finite support. Applying one primitive operator can move its support only along allowed edges. Induction on word length confines every operator image of a seed to packets reachable by a finite path. Finite sums and inverses cannot introduce a coordinate outside that set. The final assertion follows by taking the projection onto an omitted nonzero packet.
\end{proof}
Reachability is necessary, not sufficient. A permitted path may carry zero action on the particular seed, or its image may be a proper subgroup of the target packet. The theorem therefore gives a non-generation certificate but does not assert generation whenever the graph is connected.

If the packet projections are continuous into Hausdorff groups, the kernel of an omitted projection is closed. The same witness then obstructs topological generation in that topology. Later we will use precisely this argument after allowing every abstract endomorphism of the group.

\section{Observations and the original object}
An observation system also needs a contract. For a residually finite group, finite quotients separate distinct elements. This is an elementwise assertion: if two elements differ, at least one quotient distinguishes them. It does not assert that the unmarked family of quotients identifies the source isomorphism type.

The distinction is easiest to express in terms of an embedding
\[
G\longrightarrow \varprojlim G/N.
\]
Residual finiteness makes this map injective. The inverse limit includes all compatible tuples. The image of the original $G$ can nevertheless be a proper dense subgroup. Forgetting that image can forget the object even while retaining enough coordinates to distinguish every element once the object is supplied.

The later examples do not contradict the construction of an inverse limit. They show that reconstructing the complete ambient realization and identifying a specified noncomplete realization are different tasks. This is the point at which compatibility must be separated from admissibility.

\chapter{The squarefree ambient system}
\label{ch:ambient}
\chapterintro{One sequence of finite cyclic groups has two natural limits. Their generation properties differ because the connecting maps and the permitted elements differ.}

\section{Finite stages and two connecting systems}
For $N\ge1$, define
\[
A_N=\prod_{i<N}C_{p_i},\qquad m_N=\prod_{i<N}p_i.
\]
The Chinese remainder map sends $k\bmod m_N$ to $(k\bmod p_i)_{i<N}$. It is an isomorphism, so $A_N$ is cyclic and its all-one vector is a generator.

\begin{center}
\begin{tabular}{@{}cccc@{}}
\toprule
$N$&Prime factors&$m_N$&Finite group\\\midrule
1&$2$&2&$C_2$\\
2&$2,3$&6&$C_6$\\
3&$2,3,5$&30&$C_{30}$\\
4&$2,3,5,7$&210&$C_{210}$\\
5&$2,3,5,7,11$&2310&$C_{2310}$\\\bottomrule
\end{tabular}
\end{center}

Zero extension gives injections $A_N\to A_{N+1}$. Forgetting the last coordinate gives surjections $A_{N+1}\to A_N$. These are different diagrams. Their limits are
\[
D=\varinjlim_N A_N=\bigoplus_p C_p,
\qquad
P=\varprojlim_N A_N=\prod_p C_p.
\]
The first group contains only finite-support vectors. The second contains arbitrary compatible coordinate assignments.

Give each $C_p$ its discrete topology and give $P$ the product topology. A basic neighborhood fixes only finitely many coordinates. We give every subgroup of $P$ the induced topology when discussing topological generation, unless explicitly stated otherwise.

\section{Torsion and endomorphisms}
\begin{theorem}[The torsion subgroup of the product]
\label{thm:torP}
The subgroup $D$ is dense and proper in $P$, and
\[
\Tor(P)=D.
\]
\end{theorem}
\begin{proof}
A prescribed finite coordinate pattern is realized by a finite-support vector. Thus $D$ meets every basic open set and is dense. A finite-support vector is annihilated by the product of the finitely many primes in its support, so it is torsion.

Conversely, suppose $mx=0$ for a positive integer $m$. Whenever $p\nmid m$, multiplication by $m$ is invertible in $C_p$, and hence $x_p=0$. The support of $x$ is contained in the finite set of prime divisors of $m$. Thus $x\in D$. The all-one vector has infinite support, so it belongs to $P\setminus D$.
\end{proof}

\begin{proposition}[Diagonal endomorphisms of the direct sum]
\label{thm:endD}
Every endomorphism of $D$ acts by a scalar on each prime coordinate, and
\[
\End(D)\cong\prod_p\F_p
\]
as rings.
\end{proposition}
\begin{proof}
The subgroup annihilated by a prime $p$ is exactly the $C_p$ coordinate. Every endomorphism preserves it, because $px=0$ implies $pT(x)=T(px)=0$. Its restriction to that cyclic group is multiplication by a unique scalar in $\F_p$. An arbitrary element of $D$ is a finite sum of such primary elements, so these restrictions determine the map. Conversely, any family of coordinate scalars preserves finite support and defines an endomorphism. Addition and composition correspond to coordinate addition and multiplication.
\end{proof}

\begin{corollary}[An infinite seed obstruction]
\label{thm:rankD}
The endomorphism-seed rank of $D$ is $\aleph_0$. No finite set topologically generates $D$, even after all endomorphisms are admitted.
\end{corollary}
\begin{proof}
The union of the supports of finitely many seeds is finite. Diagonal endomorphisms preserve its complement as zero coordinates. Any omitted coordinate projection is a certificate from Theorem~\ref{thm:reach}. Its kernel is closed in the induced topology, so taking closure does not help. The countably many vectors $e_p$ generate $D$ algebraically, proving the exact rank.
\end{proof}
Notice that this conclusion does not depend on restricting the number of endomorphisms. Even the full endomorphism ring cannot create a missing prime component from the available seeds.

\section{A single topological generator after completion}
\begin{theorem}[Topological cyclicity]
\label{thm:procyclic}
Let $u=(1,1,1,\ldots)\in P$. Then
\[
\overline{\langle u\rangle}=P.
\]
In particular, $\trank(P)=1$.
\end{theorem}
\begin{proof}
A basic open set prescribes a tuple on finitely many distinct prime coordinates. The Chinese remainder theorem supplies an integer $k$ with the specified residues. The element $ku$ belongs to the given open set. Therefore the cyclic subgroup is dense. The product is nontrivial, so zero generators do not suffice.
\end{proof}
The generator $u$ is absent from $D$. The topological statement does not exhibit a finite generating family inside the original dense subgroup. This is the precise mechanism by which completion changes the generation contract.

\begin{figure}[tbp]
\centering\includegraphics[width=\textwidth]{\figfile{01_two_limits.pdf}}
\caption{The finite stages agree, but the arrows do not. Zero extension produces the finite-support group $D$. Coordinate restriction produces the complete product $P$. The labels state algebraic endomorphism-seed rank for $D$ and ordinary topological rank for $P$.}
\label{fig:twolimits}
\end{figure}

\section{The complete finite-quotient system}
\begin{proposition}[Profinite completion of the direct sum]
\label{thm:Dcompletion}
The canonical embedding of $D$ into $P$ identifies $P$ with $\widehat D$.
\end{proposition}
\begin{proof}
Let $f:D\to F$ be a homomorphism to a finite group. The image of a prime-order factor $C_p$ is trivial unless $p$ divides $\lvert F\rvert$. Thus $f$ vanishes on all but finitely many prime components and factors through a finite coordinate product. Kernels of finite coordinate projections are therefore cofinal among the finite-index normal subgroups of $D$. Taking their inverse limit gives exactly $P$, with the canonical map supplied by the coordinates.
\end{proof}
This proof does not assume that $F$ is abelian. The image is abelian because the source is, and the order argument applies to each factor. The definition and general properties of profinite completion may be found in \cite{Wilkes}. All needed cofinality statements in this book are proved for the specific groups under consideration.

\section{A rational space behind the torsion}
Put $V=P/D$ and let $q:P\to V$ be the quotient map.

\begin{theorem}[Unique division in the ambient quotient]
\label{thm:Vrational}
Multiplication by every positive integer is bijective on $V$. Thus $V$ has a unique rational vector-space structure extending its additive group structure. The product $P$ has no nonzero divisible subgroup.
\end{theorem}
\begin{proof}
Fix $m\ge1$ and $x\in P$. Define $y_p=m^{-1}x_p$ when $p\nmid m$, and $y_p=0$ when $p\mid m$. Then $x-my$ has finite support. Thus $q(x)=m q(y)$, proving surjectivity on $V$.

For injectivity, suppose $mx\in D$. Outside the union of the finite support of $mx$ and the prime divisors of $m$, multiplication by $m$ is invertible and $x_p=0$. Hence $x\in D$ and $q(x)=0$. Unique division defines $(a/b)v$ as the unique solution $w$ of $bw=av$. The vector-space laws follow by clearing positive denominators and using injectivity of integer multiplication.

Finally, a homomorphism from a divisible group to $C_p$ is zero. For each source element $x$, write $x=py$ and observe that its image is $p$ times an element of $C_p$. Every coordinate of a divisible subgroup of $P$ therefore vanishes, so that subgroup is zero.
\end{proof}
A group with no nonzero divisible subgroup is called \emph{reduced}. Thus $P$, and every subgroup of $P$, is reduced. Nevertheless $P/D$ is divisible. Quotients can possess division that no subgroup of the original group supports. The nonsplitting theorem later turns this distinction into an exact obstruction.

\chapter{Prime spectra and finite descriptions}
\label{ch:spectrum}
\chapterintro{The direct sum has a completely transparent subgroup universe. It separates the complexity of selecting prime components from the complexity of performing the group operation.}

\section{The entire subgroup lattice}
For $B\subseteq\N$, set
\[
D_B=\bigoplus_{i\in B}C_{p_i}\le D.
\]
A prime spectrum records which prime orders occur. Define $\sigma_p(M)=1$ when $M$ contains a nonzero element annihilated by $p$, and set it to zero otherwise. Since $p$ is prime, a nonzero such element has order exactly $p$.

\begin{theorem}[Subgroup and isomorphism classification]
\label{thm:subD}
The map $B\mapsto D_B$ is an order isomorphism from $\mathcal P(\N)$ onto $\Sub(D)$. Every subgroup is fully invariant under all endomorphisms of $D$, and
\[
D_B\cong D_C\quad\Longleftrightarrow\quad B=C.
\]
\end{theorem}
\begin{proof}
Let $H\le D$ and let $x\in H$. Its support is a finite set $F$ of primes. For $p\in F$, choose by the Chinese remainder theorem an integer $k$ that is one modulo $p$ and zero modulo every other prime in $F$. Then $kx$ is supported only at $p$, where it is nonzero. It generates the entire $C_p$ coordinate. Thus $H$ contains every coordinate factor appearing in one of its elements.

Let $B$ be the set of indices of these factors. Every element of $H$ lies in $D_B$, and every finite sum of the factors in $B$ lies in $H$. Hence $H=D_B$. Uniqueness follows by testing membership of $e_{p_i}$. Proposition~\ref{thm:endD} proves full invariance. Finally $\sigma_{p_i}(D_B)=1$ exactly when $i\in B$, and these values are isomorphism invariants.
\end{proof}
The lattice operations have literal set-theoretic descriptions:
\[
D_B+D_C=D_{B\cup C},\qquad D_B\cap D_C=D_{B\cap C},
\qquad D=D_B\oplus D_{\N\setminus B}.
\]
The infinite support choice is therefore already visible in an elementary algebraic lattice. No exotic group operation is required.

Give $\Sub(D)$ the topology inherited from $\{0,1\}^D$ by its membership functions. Membership of a vector with finite support depends on finitely many selector bits. Conversely, membership of $e_{p_i}$ recovers the $i$th bit. Consequently the order isomorphism above is also a homeomorphism from Cantor space $\bits$ to $\Sub(D)$.

\begin{corollary}[Finite observations leave continuum many possibilities]
\label{thm:finiteviews}
Fix finitely many prime indices and their membership values. There are continuum many pairwise nonisomorphic $D_B$ consistent with those values. A full compatible sequence of finite membership patterns determines one $B$, but no fixed finite prefix does.
\end{corollary}
\begin{proof}
The unspecified bits form a countably infinite set of independent binary choices. Distinct completions give nonisomorphic groups by Theorem~\ref{thm:subD}. A full compatible family fixes every coordinate and hence fixes the subset.
\end{proof}

\section{Enumeration is enough for a computable copy}
A subset of $\N$ is computably enumerable if a program can list exactly its elements. The program need not decide that an element will never appear. All these notions relativize to an oracle $X$.

\begin{theorem}[The prime-spectrum criterion]
\label{thm:spectrumce}
For every oracle $X$,
\[
D_B\text{ has an }X\text{-computable copy}
\quad\Longleftrightarrow\quad
B\text{ is computably enumerable relative to }X.
\]
\end{theorem}
\begin{proof}
Suppose first that a copy is computable in $X$. For each $i$, search its domain for a nonzero element $x$ with $p_i x=0$. Each individual test is decidable relative to $X$. Such a witness exists exactly when $i\in B$. Dovetailing the searches enumerates $B$.

Conversely, suppose $B$ is infinite and enumerable in $X$. List its distinct elements as $b_0,b_1,\ldots$. To compute $b_j$, wait for the $(j+1)$st distinct enumerated value. This procedure terminates for every $j$. Use canonical finite tuples with the $j$th coordinate in $\F_{p_{b_j}}$. Coordinatewise addition, normalization, and equality are computable in $X$. If $B$ is finite, $D_B$ is a finite group and has a computable copy. This last existence statement does not require a uniform way to decide when a finite enumeration has finished.
\end{proof}
This construction is classical. Montalb\'an gives the prime direct-sum example and its enumeration-degree interpretation explicitly \cite[Example III.39]{Montalban}. It is included here because its contrast with the later common-completion family is mathematically informative.

A noncomputable enumerable selector can therefore produce a group with computable multiplication and equality. The computable copy labels coordinates by their order of discovery. It does not supply a decision procedure for membership in the original selector.

\section{A finite definition without computable equality}
Let $K=\{n:\text{program }n\text{ halts on input }n\}$ be the diagonal halting set. It is enumerable. It is not decidable: a decision procedure would permit a program that halts on its own index precisely when the decision procedure says it does not. Its complement is not enumerable either, since simultaneous enumeration of a set and its complement decides membership.

Define
\[
B_K=\{2n:n\in\N\}\cup\{2n+1:n\notin K\},
\qquad H=D_{B_K}.
\]
This is an explicit mathematical definition of a countably infinite, locally cyclic abelian group.

\begin{corollary}
\label{thm:noncompD}
The group $H$ has no computable copy. It nevertheless has an effectively enumerable infinite list of generators and relators.
\end{corollary}
\begin{proof}
If $B_K$ were enumerable, filtering its odd entries would enumerate the complement of $K$. Theorem~\ref{thm:spectrumce} therefore excludes a computable copy.

For the relator description, introduce a commuting generator $x_i$ for each prime index, impose $p_i x_i=0$, and impose $x_{2n+1}=0$ whenever program $n$ is observed to halt on input $n$. The resulting quotient retains exactly the prime factors indexed by $B_K$. All relators are effectively enumerable. Deciding whether an odd generator has become trivial is not supplied by this enumeration.
\end{proof}
The conclusion is not absolute finite indescribability. The displayed definition is finite. What fails is the stronger contract of a countable copy with computable operations and decidable equality.

\section{A purely finite-state boundary}
Consider a restricted grammar in which a finite automaton reads the unary word $1^n$ and selects the factor $C_{p_n}$ exactly when it accepts.

\begin{theorem}[Unary selector classification]
\label{thm:unary}
A group $D_B$ is realizable by this unary-selector grammar if and only if $B$ is ultimately periodic.
\end{theorem}
\begin{proof}
On a unary input, the state after $n$ letters is obtained by iterating a single map of a finite state set. The orbit eventually repeats a state, and from that point it is periodic. Its acceptance sequence is therefore ultimately periodic.

Conversely, any finite prefix followed by a periodic binary pattern is recognized by a finite directed path entering a cycle. The prime-spectrum isomorphism criterion ensures that a different accepted set cannot accidentally realize an isomorphic $D_B$.
\end{proof}

\begin{example}[A computable group outside the unary grammar]
Let $B=\{2^k:k\ge0\}$. It is decidable and enumerable, so $D_B$ has a computable copy. It is not ultimately periodic. Indeed, if $d>0$ were an eventual period, choose $2^k$ beyond the threshold with $2^k>d$. Periodicity would force $2^k+d$ to be a power of two, although it lies strictly between $2^k$ and $2^{k+1}$. Thus $D_B$ is outside the unary grammar. A binary-input automaton recognizes the strings $10^k$, so changing the input representation removes this particular obstruction.
\end{example}

\section{What the first family does not accomplish}
The groups $D_B$ have distinct prime spectra, and those differences are already detectable by finite quotients. They do not yet show a failure of reconstruction from the full unmarked finite-quotient system. They isolate a simpler issue: which prime factors have been admitted in the first place.

The main family keeps every prime torsion factor present. It then places the selector in the divisibility of infinite-order elements instead. As a result, all finite quotients remain fixed while the effective cost of the source group varies. The rational space $P/D$ is the mechanism that permits this stronger separation.

\chapter{Pure preimages and common finite observations}
\label{ch:purity}
\chapterintro{The common-completion phenomenon is not confined to one selector construction. It follows from a general correspondence between rational subspaces and pure intermediate groups.}

\section{The preimage construction}
Keep $D\subseteq P$ and $V=P/D$ from Chapter~\ref{ch:ambient}. For a rational vector subspace $W\le V$, define
\[
H_W=q^{-1}(W).
\]
This is an actual subgroup of $P$ containing $D$. Its quotient by $D$ is canonically $W$. It is dense because it contains the dense subgroup $D$.

The construction separates the ambient coordinate space from the admissible realization. The quotient $V$ records infinite coordinate tails modulo finite changes. A rational subspace $W$ selects which such tails are allowed. The preimage restores all finite-support corrections.

\begin{definition}[Purity]
A subgroup $H$ of an abelian group $M$ is pure when
\[
H\cap mM=mH\qquad(m\ge1).
\]
Equivalently, if an equation $my=h$ with $h\in H$ has a solution in $M$, then it has a solution in $H$.
\end{definition}
Purity is a standard algebraic notion \cite{Fuchs}. The proof below uses only the displayed criterion. In particular, no structure theorem for mixed abelian groups is required.

\begin{theorem}[Pure-preimage correspondence]
\label{thm:purecorrespondence}
The map
\[
W\longmapsto H_W=q^{-1}(W)
\]
is an inclusion-preserving bijection between rational subspaces of $V$ and pure subgroups $H$ of $P$ containing $D$.
\end{theorem}
\begin{proof}
First let $W$ be rational and suppose $h\in H_W\cap mP$. Because $W$ is closed under division by $m$, choose $x\in H_W$ with $q(x)=q(h)/m$. Then $d=h-mx$ belongs to $D$. At a prime dividing $m$, the coordinate of $h$ is zero because $h\in mP$, and that of $mx$ is also zero. Thus $d_p=0$ whenever $p\mid m$.

Define $e_p=m^{-1}d_p$ at every other prime and $e_p=0$ at primes dividing $m$. The vector $e$ has finite support and satisfies $me=d$. Hence $h=m(x+e)$ with $x+e\in H_W$. This proves purity.

Conversely, let $D\subseteq H\subseteq P$ be pure. Its image $q(H)$ is an additive subgroup of $V$. Given $h\in H$ and $m\ge1$, let $d\in D$ agree with $h$ on the primes dividing $m$ and vanish elsewhere. Then $h-d\in mP\cap H$. Purity gives $h-d=mx$ for some $x\in H$, so $q(x)=q(h)/m$. Thus $q(H)$ is closed under rational scalars. Since $D\subseteq H$, the equality $H=q^{-1}(q(H))$ holds. This proves that the constructions are inverse. Both preserve inclusion.
\end{proof}
The bottom object is $H_{\{0\}}=D$, and the top object is $H_V=P$. Intersections of rational subspaces correspond to intersections of their preimages. Sums correspond to subgroup sums, with $D$ serving as the bottom of the interval of intermediate groups.

\section{Root conditions do not depend on the subspace}
\begin{theorem}[Ambient root criterion]
\label{thm:ambientroot}
For every rational subspace $W$, every $h\in H_W$, and every $m\ge1$,
\[
h\in mH_W\quad\Longleftrightarrow\quad
h_p=0\text{ for all primes }p\mid m.
\]
\end{theorem}
\begin{proof}
Multiplication by $m$ on $P$ is zero at primes dividing $m$ and invertible on every other coordinate. Thus the right-hand condition is equivalent to $h\in mP$. Apply purity from Theorem~\ref{thm:purecorrespondence}.
\end{proof}
The condition is independent of $W$, but the set of admissible $h$ is not. This distinction is the reason identical quotient laws can coexist with different original groups.

\begin{theorem}[One finite-quotient system for every preimage]
\label{thm:generalquotients}
For every rational subspace $W$ and every $m\ge1$, the coordinate map induces an isomorphism
\[
H_W/mH_W\ \cong\ \prod_{p\mid m}C_p.
\]
For $m\mid n$, the quotient transition map is the coordinate restriction from primes dividing $n$ to primes dividing $m$. These quotients are cofinal among all finite quotients, and
\[
\widehat{H_W}\cong P.
\]
\end{theorem}
\begin{proof}
The coordinate map onto the displayed finite product is surjective already on $D$. Its kernel is $mH_W$ by Theorem~\ref{thm:ambientroot}. The transition assertion follows by evaluating the same coordinates.

If $H_W\to F$ is a homomorphism to a finite group, then the image is abelian and is killed by $\lvert F\rvert$. Its kernel therefore contains $\lvert F\rvert H_W$. Hence the displayed kernels are cofinal. Their inverse limit is the product of all prime coordinates, which is $P$. Under this identification the completion map is the original inclusion $H_W\hookrightarrow P$.
\end{proof}
There is a stronger natural statement. It compares maps to every finite target rather than merely comparing the quotient groups.

\begin{proposition}[Finite-target equivalence]
\label{thm:finitehom}
For every finite group $F$, restriction along $D\hookrightarrow H_W$ gives a natural bijection
\[
\Hom(H_W,F)\longrightarrow\Hom(D,F).
\]
\end{proposition}
\begin{proof}
Put $m=\lvert F\rvert$. Every $h\in H_W$ agrees modulo $mH_W$ with an element of $D$, namely its finite coordinate vector on primes dividing $m$. Every homomorphism to $F$ kills $mH_W$. Thus two maps agreeing on $D$ agree everywhere.

Conversely, a map $D\to F$ kills every prime factor whose order does not divide $m$. It factors through the finite coordinate product on primes dividing $m$. Composing that finite-coordinate rule with the coordinate projection of $H_W$ extends it. The extension is unique by the first part. Postcomposition with a map of finite groups commutes with restriction, giving naturality.
\end{proof}
The argument never subtracts maps to a nonabelian group. The assertion is a bijection of sets of homomorphisms, and naturality is with respect to postcomposition.

\section{Countability, rank, and nonsplitting}
\begin{proposition}[Structure of an ambient preimage]
\label{thm:HWstructure}
For every rational subspace $W$,
\[
\Tor(H_W)=D,\qquad H_W/D\cong W.
\]
The group $H_W$ is reduced and residually finite. It is countable exactly when $W$ has finite or countably infinite dimension over $\Q$. If $W\ne0$, the sequence
\[
0\longrightarrow D\longrightarrow H_W\longrightarrow W\longrightarrow0
\]
does not split.
\end{proposition}
\begin{proof}
The torsion equality follows from $D\subseteq H_W\subseteq P$ and Theorem~\ref{thm:torP}. The quotient assertion is the first isomorphism theorem. Reducedness follows from reducedness of $P$. A nonzero element has a nonzero finite coordinate, which detects it and proves residual finiteness.

A rational vector space on a finite or countable basis is countable, since each vector has a finite rational linear expression. Each fiber of $H_W\to W$ is a coset of the countable group $D$. Thus $H_W$ is countable in this case. Conversely, a quotient of a countable group is countable and has at most a countable basis.

A splitting would embed the nonzero divisible group $W$ into $P$. Theorem~\ref{thm:Vrational} excludes every such subgroup.
\end{proof}
The rank of the torsion-free quotient is exactly $\dim_{\Q}W$. This gives countable examples of every finite rank and of countably infinite rank with the same finite-quotient system. The rank-one family studied next is chosen because it admits an especially sharp classification and computability theorem.

\section{What has been fixed and what still varies}
The entire product $P$, all finite quotient objects, their transition maps, and the maps to arbitrary finite targets are fixed throughout the construction. The subspace $W$ remains free. This is an ambient realization theorem, not a method for discovering $W$ from those observations.

The preimage formula also identifies the correct role of completion. Every $H_W$ contains the same dense torsion subgroup. Its completion admits every compatible coordinate tuple. Its original image admits only those tuples whose tail class lies in $W$. The missing information is an admissibility condition on the quotient, not a defect in the finite coordinates themselves.

\chapter{The paired-prime groups}
\label{ch:construction}
\chapterintro{A binary selector chooses one prime from each pair. Its indicator determines a rational line in the ambient quotient, and an explicit finite cocycle turns that line into a countable group.}

\section{Choosing a rational line}
Let $A\subseteq\N$, write $a_n=\one_A(n)$, and define
\[
S_A=\{p_{2n+a_n}:n\in\N\}.
\]
Exactly one prime from each pair is selected. Both $S_A$ and its complement are infinite. Let $u_A\in P$ be its indicator:
\[
(u_A)_p=\begin{cases}1,&p\in S_A,\\0,&p\notin S_A.\end{cases}
\]
Because $u_A$ has infinite support, $q(u_A)\ne0$. We can therefore define a rational line and its preimage by
\[
W_A=\Q q(u_A),\qquad G_A=q^{-1}(W_A).
\]
This definition produces an actual group for every set $A$. It does not claim an effective method for obtaining the selector bits of an arbitrary $A$.

\begin{figure}[tbp]
\centering\includegraphics[width=\textwidth]{\figfile{02_prime_pairs.pdf}}
\caption{Each bit selects exactly one prime from its pair. A filled point is selected. The torsion factor $C_p$ is present in $G_A$ whether or not $p$ is selected. Selection controls the infinite-order lift and its divisibility, not the presence of the torsion factor.}
\label{fig:primepairs}
\end{figure}

\section{A section with finite arithmetic defects}
For a reduced rational $r=a/b$, with $b>0$, define
\begin{equation}
\label{eq:section}
(s_A(a/b))_p=
\begin{cases}
a b^{-1}\pmod p,&p\in S_A\text{ and }p\nmid b,\\
0,&p\notin S_A\text{ or }p\mid b.
\end{cases}
\end{equation}
The inverse is used only when the denominator is a unit in $\F_p$. The convention at denominator primes is essential. It is a choice of representatives, not a claim that division by $p$ is defined in $C_p$.

The difference $b s_A(a/b)-a u_A$ is supported on the prime divisors of $b$. It follows that
\[
q(s_A(r))=r q(u_A).
\]
Thus every element of $G_A$ is $s_A(r)+d$ for some $r\in\Q$ and $d\in D$.

\begin{theorem}[Unique normal form]
\label{thm:normalform}
Every element of $G_A$ has a unique expression
\[
g=s_A(r)+d,\qquad (r,d)\in\Q\times D.
\]
Consequently $G_A$ is countably infinite, with a countable normal-form domain that does not depend on the selector.
\end{theorem}
\begin{proof}
Existence follows from the preimage definition and the displayed quotient identity. If two forms are equal, applying $q$ gives $(r-r')q(u_A)=0$. The vector $q(u_A)$ is nonzero in a rational vector space, so $r=r'$. The finite-support terms are then equal. Both $\Q$ and $D$ are countable, while the quotient is the infinite group $\Q$.
\end{proof}
The domain can be named without knowing $A$. The group operation on those names will contain the selector information.

\section{The cocycle law}
Define
\[
c_A(r,t)=s_A(r)+s_A(t)-s_A(r+t).
\]
The sum is taken in $P$, but its value belongs to $D$.

\begin{lemma}[A finite support bound]
\label{thm:cocyclesupport}
For all $r,t\in\Q$,
\[
\supp c_A(r,t)\subseteq
\{p:p\mid\den(r)\den(t)\}.
\]
\end{lemma}
\begin{proof}
At a prime not dividing either denominator, all three rationals have well-defined reductions. On a selected coordinate, reduction respects addition, so the defect is zero. On an unselected coordinate, all section values are zero. The reduced denominator of $r+t$ has no prime divisor absent from the two original denominators.
\end{proof}

\begin{theorem}[Explicit group law]
\label{thm:cocyclelaw}
In normal-form coordinates, addition is
\begin{equation}
\label{eq:grouplaw}
(r,d)+(t,e)=(r+t,d+e+c_A(r,t)).
\end{equation}
The identity is $(0,0)$ and the inverse of $(r,d)$ is $(-r,-d)$. The law is associative and commutative, and the map $(r,d)\mapsto s_A(r)+d$ is an isomorphism onto the subgroup $G_A\le P$.
\end{theorem}
\begin{proof}
Substitution into addition in $P$ gives the formula. The section satisfies $s_A(0)=0$ and $s_A(-r)=-s_A(r)$, so $c_A(r,0)=0$ and $c_A(r,-r)=0$. These identities give the stated identity and inverse.

Symmetry is immediate. The cocycle equation
\[
c_A(r,t)+c_A(r+t,v)=c_A(t,v)+c_A(r,t+v)
\]
follows by expanding and cancelling the section terms. Expanding the two associations in \eqref{eq:grouplaw} now proves associativity. Uniqueness of normal forms makes the ambient map bijective, and the addition formula makes it a homomorphism.
\end{proof}
This is the standard symmetric cocycle construction of an abelian extension \cite[Section 3.2]{Lupini}. Here the particular section makes each individual correction finitely supported and exactly computable from finitely many selector queries.

\begin{proposition}[One cocycle value reads one coordinate]
\label{thm:cocyclebit}
For every prime $p$,
\[
c_A(1/p,(p-1)/p)=-\one_{S_A}(p)e_p.
\]
\end{proposition}
\begin{proof}
The support bound leaves only the coordinate $p$. At that coordinate both fractional lifts are zero, and the lift of their sum, which is one, is $\one_{S_A}(p)$.
\end{proof}
Dropping the cocycle is therefore not a harmless simplification. The direct product operation on $\Q\times D$ removes exactly these defects and gives a different, split extension.

\section{Structure of the extension}
By Chapter~\ref{ch:purity}, the group is pure in $P$, reduced, residually finite, and dense. Its torsion is $D$, and the rational coordinate gives the exact sequence
\begin{equation}
\label{eq:extension}
0\longrightarrow D\longrightarrow G_A\xrightarrow{\rho_A}\Q\longrightarrow0,
\qquad \rho_A(s_A(r)+d)=r.
\end{equation}
It does not split.

There is an equivalent cocycle obstruction. Suppose a function $b:\Q\to D$ modified the section to an additive map $r\mapsto s_A(r)+b(r)$. Then
\[
c_A(r,t)+b(r)+b(t)-b(r+t)=0
\]
for every $r,t$. The modified section would split \eqref{eq:extension}. Hence no such function exists. The cocycle is the defect of a section into $P$, but it is not removable by a correction valued in $D$.

\begin{proposition}[Finitely generated subgroups]
\label{thm:fgsubgroups}
Every finitely generated subgroup of $G_A$ is either finite cyclic or isomorphic to $\Z\oplus C_m$ for a squarefree positive integer $m$, with $C_1$ trivial.
\end{proposition}
\begin{proof}
Let $L\le G_A$ be generated by finitely many elements. Its image in $\Q$ is a finitely generated subgroup of $\Q$, hence either zero or infinite cyclic. If it is zero, then $L\le D$ and its finite support places it inside a finite cyclic group.

Otherwise choose a lift $x\in L$ of a generator of its rational image. Subtracting suitable integer multiples of $x$ from the given generators produces finitely many elements of $D$. They generate the kernel of $L\to\rho_A(L)$, which is finite cyclic of squarefree order. The subgroup $\langle x\rangle$ intersects that torsion kernel trivially, giving the direct sum.
\end{proof}
The extensions over finitely generated rational images split, but the extension over all of $\Q$ does not. This is another exact distinction between finite algebraic stages and the global realization.

\chapter{Root certificates and finite quotients}
\label{ch:roots}
\chapterintro{An integer-root problem has a complete answer: either a finite coordinate obstructs it, or a finite correction constructs a root. This calculation controls every finite quotient.}

\begin{figure}[H]
\centering\includegraphics[width=0.94\textwidth]{\figfile{03_root_certificates.pdf}}
\caption{For the marked element $u_A$, one prime in each pair gives an obstruction and the other gives a root. A nonzero coordinate at a prime dividing the multiplier certifies impossibility. The displayed pattern is an exact worked selector, not a statistical sample.}
\label{fig:rootcertificates}
\end{figure}

\section{An exact decision criterion}
The general purity argument already establishes the root condition. For $G_A$, normal forms make the successful branch constructive.

\begin{theorem}[Root or obstruction]
\label{thm:rootcertificate}
Let $g=s_A(r)+d\in G_A$ and let $m\ge1$. Then
\[
g\in mG_A\quad\Longleftrightarrow\quad
 g_p=0\quad\text{for every }p\mid m.
\]
If the condition fails, a prime $p\mid m$ with $g_p\ne0$ is a nonexistence certificate. If it holds, a root is obtained by the finite procedure in the proof.
\end{theorem}
\begin{proof}
Necessity follows because multiplication by $m$ is zero in every coordinate $C_p$ with $p\mid m$.

Assume the condition holds and put $t=r/m$. The vector
\[
h=g-ms_A(t)
\]
has zero image in $V$ and therefore lies in $D$. Its support is contained in the union of the support of $d$ and the prime divisors of $\den(r)\den(t)$. Indeed, at any other selected coordinate, reduction gives $r-mt=0$, and unselected section values are zero.

For $p\mid m$, the coordinate $h_p$ is zero. At every other prime put $e_p=m^{-1}h_p$, and at primes dividing $m$ put $e_p=0$. Then $e\in D$ and $me=h$. Consequently
\[
y=s_A(t)+e\quad\text{satisfies}\quad my=g.
\]
Every nonzero correction coordinate is found in the stated finite support set.
\end{proof}
A failed root is witnessed by one coordinate. A successful root requires a rational division and a finite-support adjustment. This is a simple instance of a realization problem with both positive and negative certificates.

\begin{proposition}[All roots and their number]
\label{thm:rootfiber}
If $my=g$ is solvable, its solution set is the coset
\[
y_0+G_A[m],\qquad
G_A[m]=\bigoplus_{p\mid m}C_p.
\]
It has exactly $\rad(m)$ elements.
\end{proposition}
\begin{proof}
Two solutions differ by an element annihilated by $m$, and adding such an element to a solution gives another. The ambient product has a nonzero $m$-torsion coordinate only when $p\mid m$. All those finite coordinate factors lie in $D\subseteq G_A$. The kernel is therefore the displayed finite direct sum, whose cardinal is $\rad(m)$.
\end{proof}
Repeated prime powers do not enlarge the root kernel. For example, a solvable fourth-root equation has two solutions, not four. The squarefree coordinate structure is responsible for this behavior.

\section{The quotient formula}
\begin{theorem}[All integer quotients]
\label{thm:quotientformula}
For $m\ge1$, the coordinate map
\[
\pi_m:G_A\longrightarrow\prod_{p\mid m}C_p
\]
is surjective with kernel $mG_A$. Thus
\[
G_A/mG_A\cong\prod_{p\mid m}C_p\cong C_{\rad(m)}.
\]
For $m\mid n$, the transition map is the corresponding coordinate projection.
\end{theorem}
\begin{proof}
A vector supported on the prime divisors of $m$ realizes any tuple in the target, proving surjectivity. The root criterion identifies the kernel. The Chinese remainder theorem identifies the finite product with the squarefree cyclic group. The transition assertion follows by taking the same coordinates before and after quotienting.
\end{proof}

\begin{center}
\begin{tabular}{@{}rrrr@{}}
\toprule
$m$&$\rad(m)$&$G_A/mG_A$&Root-fiber size when nonempty\\\midrule
1&1&$0$&1\\
4&2&$C_2$&2\\
8&2&$C_2$&2\\
12&6&$C_6$&6\\
18&6&$C_6$&6\\
30&30&$C_{30}$&30\\
60&30&$C_{30}$&30\\\bottomrule
\end{tabular}
\end{center}

The kernel itself depends only on the distinct prime divisors:
\[
mG_A=\rad(m)G_A.
\]
This is stronger than equality of the quotient orders. It is equality of the subgroups inside the fixed source group.

\section{Why these are all finite quotients}
Let $f:G_A\to F$ be any homomorphism to a finite group. Its image is an abelian subgroup of $F$, and every element of the image is annihilated by $\lvert F\rvert$. Therefore $f$ kills $\lvert F\rvert G_A$ and factors through the quotient in Theorem~\ref{thm:quotientformula}.

This proves that every finite quotient is cyclic of squarefree order. Conversely every cyclic squarefree group occurs, by selecting its finite prime set. The quotient system, not merely this list, is canonically the same for all selectors.

\begin{corollary}[Common profinite completion]
\label{thm:GAcompletion}
For every $A\subseteq\N$, the inclusion $G_A\subseteq P$ extends to a canonical isomorphism
\[
\widehat{G_A}\cong P.
\]
The induced topology on $G_A$ is its full profinite topology.
\end{corollary}
\begin{proof}
The coordinate kernels are cofinal among finite-index normal subgroups by the preceding factorization argument. Their inverse limit is $P$. Neighborhoods defined by finitely many coordinate zeros coincide with the kernels $mG_A$, so the induced and profinite topologies agree.
\end{proof}
The product $P$ has a computable profinite presentation: list the primes, form the finite addition tables of $A_N$, and give the coordinate projections. This is the standard inverse-system meaning of a computable profinite presentation \cite[Definition 3.3]{Lupini}. It is independent of $A$. The uncountable product is presented by an effective inverse system, not by a computable listing of all its elements. This does not turn an arbitrary $G_A$ into a computable discrete group.

\section{The information retained by a marked element}
For the distinguished element $u_A=s_A(1)$, the coordinate criterion simplifies to
\begin{equation}
\label{eq:markroot}
p\mid_{G_A}u_A\quad\Longleftrightarrow\quad p\notin S_A.
\end{equation}
Consequently, for each $n$,
\begin{align}
 p_{2n}\mid_{G_A}u_A&\quad\Longleftrightarrow\quad n\in A,\label{eq:evenroot}\\
 p_{2n+1}\mid_{G_A}u_A&\quad\Longleftrightarrow\quad n\notin A.\label{eq:oddroot}
\end{align}
Exactly one root equation is solvable. This binary pattern is the selector written in intrinsic group language.


The finite quotient objects alone have not changed. What has been added is the coherent image of a particular source element. It is therefore important not to paraphrase the common-completion theorem as equality of every possible finite observation. Marked source data is a stronger observation contract.

\chapter{Classification, morphisms, and seed rank}
\label{ch:classification}
\chapterintro{The finite quotients are fixed, but the original groups are not. Divisibility classifies their isomorphism types, and the coordinate structure determines every homomorphism.}

\section{The signature of an infinite-order element}
The rational coordinate of a normal form detects torsion: $s_A(r)+d$ has infinite order exactly when $r\ne0$.

\begin{lemma}[Tail signature]
\label{thm:tailsignature}
If $g=s_A(r)+d$ has infinite order, then for all but finitely many primes,
\[
p\mid_{G_A}g\quad\Longleftrightarrow\quad p\notin S_A.
\]
An explicit exception set consists of the primes dividing the nonzero numerator or denominator of $r$, together with $\supp(d)$.
\end{lemma}
\begin{proof}
Outside that set, the reduction of $r$ is defined and nonzero. The coordinate of $g$ is therefore nonzero exactly at the selected primes. Apply Theorem~\ref{thm:rootcertificate} with multiplier $p$.
\end{proof}
The full signature of the distinguished element is exact. The signature of an arbitrary infinite-order element is exact only in the tail. This finite discrepancy is what determines the correct unmarked classification.

\section{Isomorphism means eventual agreement}
\begin{theorem}[Unmarked isomorphism classification]
\label{thm:isoclass}
For all $A,B\subseteq\N$,
\[
G_A\cong G_B\quad\Longleftrightarrow\quad A\eqfin B.
\]
When $A\eqfin B$, the two groups are equal as subgroups of the fixed product $P$.
\end{theorem}
\begin{proof}
Suppose $f:G_A\to G_B$ is an isomorphism. The image of $u_A$ has infinite order. An isomorphism preserves the existence of a $p$th root for every prime $p$. Combining the exact signature \eqref{eq:markroot} with Lemma~\ref{thm:tailsignature} for $f(u_A)$ gives $S_A\eqfin S_B$. Since each set chooses exactly one prime from each pair, the corresponding sets of bits differ finitely.

Conversely, suppose $A\eqfin B$. For every rational $r$, the vectors $s_A(r)$ and $s_B(r)$ differ only on the finitely many prime pairs where the selectors differ. Thus their difference belongs to $D$. The descriptions $G_A=s_A(\Q)+D$ and $G_B=s_B(\Q)+D$ define the same subset of $P$.
\end{proof}
The normal forms need not be identical. Their coordinate change is
\begin{equation}
\label{eq:transport}
(r,d)\longmapsto
\bigl(r,d+s_A(r)-s_B(r)\bigr).
\end{equation}
It fixes the represented ambient element. The correction is finite exactly because the selectors differ only finitely.

\begin{corollary}[Marked classification]
\label{thm:markedclass}
An isomorphism $(G_A,u_A)\cong(G_B,u_B)$ exists if and only if $A=B$.
\end{corollary}
\begin{proof}
A marked isomorphism preserves the exact two-prime root test \eqref{eq:evenroot}--\eqref{eq:oddroot} at every index. Therefore every bit agrees. The converse is the identity map.
\end{proof}

\begin{example}[The same group with a different marked point]
Let $A=\varnothing$ and $B=\{0\}$. Only the first prime pair changes. Then $G_A=G_B$ inside $P$, but $u_A-u_B=e_2-e_3$. The distinguished element is not preserved. In particular, the statement $G_A\cong G_B\iff A=B$ would be false for unmarked groups.
\end{example}

\section{All homomorphisms are diagonal}
The product $P$ is also a commutative ring under coordinatewise multiplication. Multiplication by any $\lambda\in P$ preserves $D$, so it induces an additive endomorphism $\overline M_\lambda$ of $V$. Because division in $V$ is unique, this induced endomorphism is $\Q$-linear.

\begin{theorem}[Complete homomorphism description]
\label{thm:allhom}
Every homomorphism $f:G_A\to G_B$ has the form
\[
f(g)_p=\lambda_p g_p\qquad(p\text{ prime})
\]
for a unique $\lambda\in P$. Conversely, coordinate multiplication by $\lambda$ defines such a homomorphism exactly when
\begin{equation}
\label{eq:homcondition}
q(\lambda u_A)\in\Q q(u_B).
\end{equation}
Equivalently, there is a rational $r$ for which $\lambda u_A-s_B(r)$ has finite support.
\end{theorem}
\begin{proof}
For a homomorphism $f$, the composite $G_A\to G_B\to C_p$ kills $pG_A$. The quotient $G_A/pG_A$ is canonically $C_p$. Every endomorphism of this cyclic additive group is multiplication by a unique scalar $\lambda_p$. Hence the displayed coordinate rule holds. Uniqueness also follows by evaluating $f(e_p)$.

Conversely, multiplication by any $\lambda$ is a group endomorphism of $P$ and preserves $D$. Since its induced map on $V$ is rational-linear, it carries $\Q q(u_A)$ into $\Q q(u_B)$ exactly when it carries the generating vector into that line. This is \eqref{eq:homcondition}. Taking preimages proves that the restriction maps $G_A$ into $G_B$. The equivalent finite-support condition follows from the definition of the quotient.
\end{proof}
No continuity assumption was imposed. Every resulting homomorphism is nevertheless continuous for the induced profinite topologies, because each output coordinate depends only on the corresponding input coordinate. Its extension to the completion is the same coordinatewise scalar map on $P$.

\begin{theorem}[Torsion-image obstruction]
\label{thm:homobstruction}
If $A\neqfin B$, every homomorphism $G_A\to G_B$ has image contained in $D$. Consequently neither group embeds in the other, and neither is a quotient of the other.
\end{theorem}
\begin{proof}
A homomorphism preserves torsion and induces a map $\Q\to\Q$ on the quotients, necessarily multiplication by a rational $r$. If $r=0$, its image lies in $D$.

Assume $r\ne0$. Then $f(u_A)$ has the tail signature of $S_B$. For every $p\notin S_A$, the element $u_A$ has a $p$th root, and its image has a $p$th root as well. Lemma~\ref{thm:tailsignature} implies that $S_B\setminus S_A$ is finite. Each differing prime pair contributes one element to this difference, so $A\eqfin B$, contrary to the assumption.

An embedding cannot send an infinite-order element to torsion. A surjection onto $G_B$ cannot have torsion image, since $G_B/D\cong\Q\ne0$.
\end{proof}
A useful precise consequence of Theorem~\ref{thm:allhom} is that, when $A\neqfin B$, the possible diagonal sequences are exactly those for which $\lambda_p=0$ at all but finitely many primes of $S_A$. Their values on the complementary primes are unrestricted. These maps can be nonzero, and their torsion images need not be finite.

\section{The endomorphism ring and its units}
Write $\overline r_p$ for reduction of a rational $r$ at a prime not dividing its denominator. The phrase ``for almost every selected prime'' below means outside a finite set, which may include those denominator primes.

\begin{corollary}[Endomorphism ring]
\label{thm:endGA}
Under coordinate multiplication,
\[
\End(G_A)\cong
E_A=\{\lambda\in P:\exists r\in\Q,\ 
\lambda_p=\overline r_p\text{ for almost every }p\in S_A\}.
\]
The rational $r$ is unique. It is the scalar induced on $G_A/D$, and the map $E_A\to\Q$ sending $\lambda$ to $r$ is a surjective ring homomorphism.
\end{corollary}
\begin{proof}
Apply Theorem~\ref{thm:allhom} with $B=A$. At unselected coordinates, both $\lambda u_A$ and $s_A(r)$ vanish. At selected coordinates the finite-support condition is exactly the displayed eventual equality.

If two different rationals agreed modulo almost every selected prime, their nonzero difference would have a numerator divisible by infinitely many primes, which is impossible. Thus $r$ is unique. Addition and composition induce addition and multiplication of the quotient scalars, giving a ring homomorphism. The diagonal sequence $s_A(r)$ realizes every rational $r$, proving surjectivity.
\end{proof}
The arbitrary coordinates outside $S_A$ show why the endomorphism ring may contain infinitely specified operators. The seed-rank obstruction below remains true even when all these operators are available.

\begin{corollary}[Automorphisms]
\label{thm:autGA}
An endomorphism with diagonal $\lambda\in E_A$ is an automorphism if and only if every coordinate $\lambda_p$ is nonzero. In that case its rational quotient scalar is nonzero, and its inverse has diagonal $\lambda_p^{-1}$.
\end{corollary}
\begin{proof}
An automorphism restricts to an automorphism of each torsion factor, so every scalar is nonzero. Conversely, if every scalar is nonzero, the eventual rational scalar cannot be zero because $S_A$ is infinite. Outside finitely many selected primes, the inverse diagonal is reduction of its reciprocal. Thus the inverse diagonal belongs to $E_A$ and defines an inverse endomorphism.
\end{proof}

\section{Infinite seed rank before completion}
\begin{theorem}[Seed-rank separation]
\label{thm:rankGA}
For every $A$,
\[
\srank_{\End(G_A)}(G_A)=\aleph_0.
\]
No finite family topologically generates $G_A$, even after arbitrary endomorphisms are admitted. Yet its completion has ordinary topological generating rank one.
\end{theorem}
\begin{proof}
Each element $s_A(r)+d$ has only finitely many nonzero coordinates outside $S_A$. A finite family therefore omits infinitely many unselected prime coordinates. By Theorem~\ref{thm:allhom}, every endomorphism preserves a zero coordinate as zero. Finite sums of orbit elements still vanish at any omitted prime $p$. The projection onto $C_p$ is a non-generation certificate, since $e_p\in G_A$ has nonzero projection. Its kernel is closed, so taking closure does not repair the failure.

Countably many seeds suffice because $G_A$ is countable. For ordinary topological generation, the coordinate vectors already generate the dense subgroup $D$. Finally, the completion is $P$, whose all-one vector is a topological generator by Theorem~\ref{thm:procyclic}.
\end{proof}
The same infinite seed rank occurs for every selector. It does not measure the varying effective complexity. Algebraic generation, finite description, and computable presentation remain separate invariants.

\chapter{The exact cost of an effective realization}
\label{ch:computability}
\chapterintro{The selector does not change the completion, but it determines exactly which oracles can compute the original group. The proof must distinguish a named element from an unmarked isomorphism type.}

\section{Oracles, copies, and degrees}
An oracle for a set $X\subseteq\N$ answers membership questions about $X$. The notation $A\le_T X$ means that a terminating oracle procedure decides membership in $A$. Two sets have the same Turing degree when each is computable from the other.

The degree spectrum of a countable structure is the set of degrees of its copies. When that spectrum has a least degree, it is called the degree of the isomorphism type. This terminology is standard in computable structure theory \cite{Montalban,Calvert}. Here we prove an exact oracle criterion that specifies which oracles compute at least one copy. We do not identify this statement with an unproved claim that every such oracle degree is realized as the exact degree of a particular copy.

For a group on a countable coded domain, knowledge of its atomic diagram is equivalent, at the level needed here, to knowledge of its group operation and equality. A computation of integer multiples uses finitely many additions. A search for a root can therefore be performed in any oracle-computable copy.

\section{A presentation from the selector}
Reduced rationals have canonical integer codes. Finite-support vectors in $D$ have canonical codes given by increasing prime labels and nonzero residues. Their Cartesian product is a computably named countable domain with decidable equality and a computable bijection to $\N$.

\begin{proposition}[Uniform construction relative to the selector]
\label{thm:computeupper}
If $A\le_T X$, then the normal-form construction gives an $X$-computable copy of $G_A$. In its canonical normal-form presentation, the group operation has Turing degree exactly $\deg_T(A)$.
\end{proposition}
\begin{proof}
To add $(r,d)$ and $(t,e)$, compute the prime divisors of $\den(r)\den(t)$. At each of these finitely many primes, query the corresponding selector bit and evaluate the three section residues. Lemma~\ref{thm:cocyclesupport} guarantees that every other correction coordinate is zero. Formula~\eqref{eq:grouplaw} then gives the sum. Negation and equality are computable without querying $A$.

For the lower bound on this particular presentation, compute the sum of the coded elements $(1/p,0)$ and $((p-1)/p,0)$. Its finite torsion part is the value in Proposition~\ref{thm:cocyclebit}. For $p=p_{2n}$, that value is zero exactly when $n\in A$. Thus the operation computes every selector bit.
\end{proof}
The number of oracle questions used by one addition is finite and bounded by the number of distinct prime divisors of the denominator product. This is a locality statement about the operation, not a claim of a polynomial-time implementation. Factoring integers and finding their prime indices contribute ordinary computational costs.

\section{Decoding a marked copy}
\begin{proposition}[Uniform marked decoding]
\label{thm:markeddecode}
Given an $X$-computable copy of $G_A$ together with the name of $u_A$, one can uniformly compute $A$ in $X$.
\end{proposition}
\begin{proof}
For input $n$, search the domain simultaneously for solutions of
\[
p_{2n}y=u_A,\qquad p_{2n+1}z=u_A.
\]
Each candidate test is decidable in $X$. Exactly one equation has a solution by \eqref{eq:evenroot}--\eqref{eq:oddroot}. Dovetailing the searches therefore terminates. A solution to the first equation means $n\in A$, and a solution to the second means $n\notin A$.
\end{proof}
The procedure does not decide a negative existential statement directly. It uses two complementary positive searches. The paired selector is what converts enumerable root information into a total decision procedure.

\begin{figure}[tbp]
\centering\includegraphics[width=\textwidth]{\figfile{04_degree_equivalence.pdf}}
\caption{The two effective directions. Selector access computes the finite cocycle and hence a normal-form copy. A marked copy computes the selector by two complementary root searches. The middle-to-right arrow records an abstract isomorphism, not an algorithm for finding one. The unmarked converse uses a fixed infinite-order element and finite exception data.}
\label{fig:degrees}
\end{figure}

\section{Decoding an unmarked copy}
\begin{theorem}[Exact oracle criterion]
\label{thm:degree}
For every $A,X\subseteq\N$,
\[
G_A\text{ has an }X\text{-computable copy}
\quad\Longleftrightarrow\quad A\le_T X.
\]
In particular, the isomorphism type of $G_A$ has least presentation degree $\deg_T(A)$.
\end{theorem}
\begin{proof}
The reverse implication is Proposition~\ref{thm:computeupper}. For the forward implication, let $H$ be an $X$-computable copy. Choose the name $h$ of an infinite-order element. Such an element exists because $H$ has quotient $\Q$ by its torsion subgroup.

Under an isomorphism to $G_A$, the element $h$ has a normal form with nonzero rational part. Lemma~\ref{thm:tailsignature} gives a finite exception set. Choose $N$ so that every prime in the pairs indexed by $n\ge N$ avoids that set. For these $n$, exactly one of
\[
p_{2n}y=h,\qquad p_{2n+1}z=h
\]
has a solution, and the first is solvable exactly when $n\in A$.

Dovetailing the two searches decides every bit beyond $N$. Hard-code the name of $h$, the number $N$, and the finitely many bits below $N$ into the oracle procedure. These are finite constants, so the result is a Turing reduction of $A$ to $X$.

The canonical normal-form copy has exact degree $\deg_T(A)$ by Proposition~\ref{thm:computeupper}, while every copy computes $A$ by the forward implication. Hence this degree is the least one.
\end{proof}
There is no claim of a uniform algorithm finding an infinite-order element or its exception bound in every unmarked copy. The theorem is an existence statement about Turing reductions, and a finite list of constants is legitimate in such a reduction. For marked copies, Proposition~\ref{thm:markeddecode} gives the genuinely uniform version.

The distinction also matches the algebra. If $A\eqfin B$, then the unmarked groups are identical inside $P$ and $A$ and $B$ have the same Turing degree. Recovering a literal selector uniformly from the unmarked group would demand more than the isomorphism type remembers.

\section{Every degree over one computable completion}
\begin{corollary}[Degree variation with fixed observations]
\label{thm:alldegrees}
Every Turing degree occurs as the least presentation degree of a countable, reduced, residually finite abelian group whose profinite completion is $P$, equipped with the fixed computable inverse-system presentation of Chapter~\ref{ch:roots}.
\end{corollary}
\begin{proof}
Choose any set $A$ of the desired degree and apply Theorem~\ref{thm:degree}. All the structural and completion properties were proved in Chapters~\ref{ch:purity}--\ref{ch:roots}.
\end{proof}
For $A=\varnothing$, the copy is computable. For $A=K$, the group has no computable copy and its least degree is the halting degree $\mathbf 0'$. Iterated jumps yield further examples. An arbitrary selector supplies any desired degree without changing the finite quotient tables or their transition maps.

The group $G_K$ still has a finite mathematical definition. Its noncomputability is not a proof that it cannot be named in a finite semantic language. The theorem concerns effective copies with computable operations and equality.

\section{A sharper separation than prime-spectrum coding}
For the initial groups $D_B$, an enumeration of the selected prime set suffices. For $G_A$, the paired complementary tests decide the selector. Meanwhile, every prime torsion factor is present in every $G_A$, so the information is no longer stored in which finite-order elements exist.

It resides in the relationship between torsion coordinates and infinite-order elements. All finite quotient groups remain fixed, but the root pattern of an admissible marked element distinguishes the selectors. This is why the construction measures a realization cost rather than merely relabeling the finite quotient data.

\chapter{What finite observations cannot reconstruct}
\label{ch:reconstruction}
\chapterintro{The obstruction is not merely computational. The observation system itself has forgotten which dense realization was intended. No decoder using only that system can restore the missing choice.}

\section{The observation contract}
Let $\Obs(G)$ denote the unmarked inverse system of finite quotient groups and their transition maps, considered up to isomorphism as a pro-object. The source group $G$ is not supplied again as part of the input. Nor do we retain evaluations of named source elements. Supplying the source itself would make a reconstruction problem vacuous, and supplying coherent marked elements is a stronger contract.

For the family of this book, there is no ambiguity about the comparison of pro-objects. The same concrete finite products
\[
\prod_{p\mid m}C_p
\]
with the same coordinate projections form a cofinal model for every $\Obs(G_A)$. Proposition~\ref{thm:finitehom} also identifies the homomorphism sets into every finite target by a natural bijection induced by restriction to $D$.

\begin{theorem}[No decoder from the unmarked finite-quotient system]
\label{thm:nodecoder}
There is no isomorphism-invariant assignment $R$ on these observation systems such that
\[
R(\Obs(G_A))\cong G_A
\qquad\text{for every }A\subseteq\N.
\]
This remains impossible without any computability restriction on $R$.
\end{theorem}
\begin{proof}
Take $A=\varnothing$ and $B=\N$. Their observation systems are isomorphic by Theorem~\ref{thm:quotientformula}. Their selectors differ infinitely, so their groups are not isomorphic by Theorem~\ref{thm:isoclass}. An isomorphism-invariant assignment has one output isomorphism type on that input observation type. It cannot recover both nonisomorphic source types.
\end{proof}
The proof uses two objects, not a counting argument. The full family shows that the ambiguity is much larger: one observation system supports continuum many isomorphism types, including examples with every possible least presentation degree.

\begin{figure}[tbp]
\centering\includegraphics[width=\textwidth]{\figfile{05_shared_completion.pdf}}
\caption{Different selectors, including computable and noncomputable ones, lead to the same complete finite-quotient system and the same product completion. The arrows are completion maps. They are not isomorphisms from the countable source groups to the uncountable product.}
\label{fig:sharedcompletion}
\end{figure}

\section{Element separation does not identify the source}
Every $G_A$ is residually finite. Its coordinate maps distinguish any two distinct elements. They give an injective map
\[
G_A\hookrightarrow\varprojlim_N A_N=P.
\]
There is no conflict with Theorem~\ref{thm:nodecoder}. The injective map identifies each element once the domain is known. The unmarked inverse system does not identify the image of that domain among all compatible tuples.

The complete product is uncountable. Each $G_A$ is countable and dense, hence proper and noncomplete. In particular, a theorem asserting that a complete compatible realization equals its inverse limit has not failed. These examples deliberately retain a different admissible subspace before completion.

\begin{proposition}[A proper map invisible after completion]
\label{thm:propercompletion}
The inclusion $D\hookrightarrow G_A$ is not an isomorphism, but it induces an isomorphism of profinite completions.
\end{proposition}
\begin{proof}
The quotient $G_A/D\cong\Q$ is nonzero, so the inclusion is proper. Both completions are canonically $P$, and the coordinate inclusions agree, so the induced completion map is the identity on $P$.
\end{proof}
This supplies a map-level failure: completion need not reflect isomorphisms. It is stronger than merely exhibiting two abstract groups with isomorphic completions.

\section{Faithful on maps, but not a full reconstruction}
There is a useful categorical distinction. A functor is faithful when different maps remain different. This does not require different objects to remain nonisomorphic, and it does not require every map between observed objects to lift.

\begin{proposition}[Faithfulness and failure of fullness]
\label{thm:faithful}
Profinite completion is faithful on the category of residually finite groups and group homomorphisms. Its restriction to the full subcategory with objects $G_A$ is not full. On the larger full subcategory with objects $D$ and all $G_A$, it does not reflect isomorphisms.
\end{proposition}
\begin{proof}
Let $f,g:H\to K$ be distinct homomorphisms between residually finite groups. Choose $h$ with $f(h)\ne g(h)$. Some finite quotient of $K$ distinguishes those elements. Therefore the induced completion maps differ on the image of $h$, proving faithfulness.

For failure of fullness, choose $A\neqfin B$. Under the canonical identifications of the completions with $P$, consider the identity map $P\to P$. If it were induced by a homomorphism $G_A\to G_B$, that homomorphism would be the identity on every coordinate and hence would include $G_A$ into $G_B$. Theorem~\ref{thm:homobstruction} excludes such an embedding. After adjoining $D$ to the source category, the proper inclusion of Proposition~\ref{thm:propercompletion} proves failure to reflect isomorphisms.
\end{proof}
The functorial extension of a homomorphism to profinite completions is standard \cite[Proposition 3.1.2]{Wilkes}. The distinction between the two source categories is essential. Within the paired-prime family alone, a supplied map can be recovered more strongly.

\begin{proposition}[Reflection within the paired-prime family]
\label{thm:rankonereflection}
For a homomorphism $f:G_A\to G_B$, the induced completion map $\widehat f$ is an isomorphism if and only if $f$ is an isomorphism.
\end{proposition}
\begin{proof}
One implication follows from functoriality. Conversely, suppose $\widehat f$ is an isomorphism. By Theorem~\ref{thm:allhom}, its coordinate action is multiplication by a diagonal sequence $\lambda$. Each $\lambda_p$ must be nonzero, so $f$ is injective. If $A\neqfin B$, Theorem~\ref{thm:homobstruction} would force its image to be torsion, contradicting injectivity on an infinite-order element. Hence $A\eqfin B$, and the source and target are equal as subgroups of $P$. The map is now an endomorphism of that group with every diagonal entry nonzero. Corollary~\ref{thm:autGA} makes it an automorphism.
\end{proof}
Thus completion forgets distinctions between the objects $G_A$, and some maps between their completions have no lift. But when a homomorphism between two members of this family is already supplied, an isomorphism after completion certifies that the supplied map was an isomorphism. This does not reconstruct an unknown source from its completion alone.

\section{What a marked coherent point restores}
The coherent image of $u_A$ in the finite quotients records, at prime $p$, whether its coordinate is zero or one. Within the declared family, these zero patterns recover $S_A$ and therefore $A$.

This observation should be read in two directions. A marked coherent point can restore information discarded by the unmarked system. But specifying that infinite coherent point can itself require a noncomputable selector. The finite quotient tables remain computable while the chosen point does not.

One point is sufficient to identify our family because the construction contract says that the rational line through its tail class is the admissible quotient. Without that contract, a point in $P$ would not identify an arbitrary dense subgroup containing it. The realization rule and its parameter must be retained together.

\section{Two different inverse-limit assertions}
For the subgroup family $D_B$, a complete coherent sequence of finite membership decisions determines $B$ exactly. For the family $G_A$, the unmarked finite quotient system is independent of $A$. These are different observations, even though both involve finite structures and inverse limits.

The first system remembers which coordinate factors were selected. The second remembers the finite quotient objects but forgets how a distinguished infinite-order source element sits in them. Treating these as the same notion of ``all finite information'' would erase the very distinction that the construction proves.

\begin{keybox}
\textbf{Reconstruction boundary.} Compatibility identifies the complete ambient object. Reconstruction of a specified noncomplete source additionally requires its admissible image. A finite or effective realization must account for the information in that image.
\end{keybox}

\chapter{Escaping a fixed finite grammar}
\label{ch:grammar}
\chapterintro{A relative diagonal argument can be carried out entirely inside one common-completion class. Because finite selector changes leave the group unchanged, the diagonal must disagree infinitely often.}

\section{Fix the semantics before diagonalizing}
Let $\mathfrak G$ be a grammar with countably many finite, parameter-free descriptions. Each valid description has one group isomorphism type as its semantics. Enumerate the strings as
\[
d_0,d_1,d_2,\ldots
\]
and write $R_j$ for the semantic output. Invalid strings may be assigned a fixed dummy group. No decidability assumption is made about validity or evaluation.

For each $j$ whose output lies in the family, choose $B_j\subseteq\N$ such that $R_j\cong G_{B_j}$. For other $j$, set $B_j=\varnothing$. This is a set-theoretic choice of representatives. It is not a procedure for deciding whether an arbitrary description denotes a group in the family.

The isomorphism criterion says that a selector must avoid the whole eventual-equivalence class of $B_j$, not merely disagree at one coordinate. A naive Cantor diagonal with one disagreement for each $j$ would not establish that requirement.

\section{Repeated diagonal coding}
Fix a computable pairing bijection $\langle j,k\rangle:\N^2\to\N$. For a freely chosen set $C\subseteq\N$, define a new selector $A_C$ by
\begin{align}
\one_{A_C}(2\langle j,k\rangle)
 &=1-\one_{B_j}(2\langle j,k\rangle),\label{eq:diageven}\\
\one_{A_C}(2\langle j,k\rangle+1)
 &=\one_C(j).\label{eq:diagodd}
\end{align}
Every natural number has exactly one of these forms, so the rules define a total subset. The even positions defeat grammar outputs. The odd positions retain an independent continuum of choices, with each chosen bit repeated infinitely often.

\begin{theorem}[Escape within one profinite-completion class]
\label{thm:escape}
For every $C\subseteq\N$, the group $G_{A_C}$ is not isomorphic to any output of $\mathfrak G$. Distinct $C$ give pairwise nonisomorphic groups, with no embeddings in either direction. Every group is countable, reduced, residually finite, and has profinite completion $P$.
\end{theorem}
\begin{proof}
Fix a description index $j$. If $R_j$ is outside the family, it cannot be isomorphic to $G_{A_C}$. Otherwise rule \eqref{eq:diageven} makes $A_C$ disagree with $B_j$ at every index $2\langle j,k\rangle$. These are infinitely many distinct indices, so $A_C\neqfin B_j$. Theorem~\ref{thm:isoclass} excludes the isomorphism.

If $C\ne C'$, choose $j$ at which they differ. Rule \eqref{eq:diagodd} makes $A_C$ and $A_{C'}$ differ at infinitely many odd indices. Theorems~\ref{thm:isoclass} and \ref{thm:homobstruction} then give nonisomorphism and nonembeddability. The remaining properties hold for every paired-prime group.
\end{proof}

\begin{figure}[tbp]
\centering\includegraphics[width=\textwidth]{\figfile{06_repeated_diagonal.pdf}}
\caption{Two infinite coordinate blocks are allocated to every $j$. On the even block the selector opposes the representative $B_j$ forever. On the odd block it repeats the free bit $C(j)$ forever. Both repetitions are necessary to survive the quotient by finite changes.}
\label{fig:diagonal}
\end{figure}

\section{The complexity consequence}
The theorem implies
\[
K_{\mathfrak G}(G_{A_C})=+\infty
\]
for every $C$. It produces continuum many distinct isomorphism types outside the grammar while preserving all the finite quotient data.

This is a statement relative to a specified language and semantics. The definition of $A_C$ may use the entire semantic table of that language in the surrounding mathematical theory. It does not yield an effective evaluator for the table. Nor does it produce a paradox by defining in one paragraph an object claimed to be beyond every possible finite definition. The object escapes the fixed internal grammar, not all larger metatheories.

The same distinction applies to a language extended by a new constant. A constant naming one previously absent group enlarges the language. If it is only a semantic constant, it need not give algorithms for that group's operations. A new realization rule must be assessed under its actual semantic and effective contract.

\section{Monotonicity under grammar extension}
\begin{proposition}[Admissible expansion]
\label{thm:grammarextension}
Suppose every valid description in $\mathfrak G$ translates to a valid description in $\mathfrak H$ with the same semantic group. Then every $\mathfrak G$-realizable group is $\mathfrak H$-realizable. If every translated description has length at most $f(\lvert d\rvert)$, then
\[
K_{\mathfrak H}(G)\le f(K_{\mathfrak G}(G))
\]
whenever $K_{\mathfrak G}(G)$ is finite.
\end{proposition}
\begin{proof}
Translate a shortest valid $\mathfrak G$ description of $G$. Its semantic output remains $G$, and its length has the stated bound.
\end{proof}
An obstruction to a small grammar can therefore disappear after an admissible expansion. The power-of-two selector from Chapter~\ref{ch:spectrum} illustrates this concretely: a unary finite-state selector fails, while a binary one succeeds. Such a result characterizes a boundary rather than discrediting the larger realization program.

\section{Uniform families and individual instances}
The expression $A\mapsto G_A$ is one uniform mathematical construction. It is not a finite description of each arbitrary instance unless the selector is also described. A single family can contain uncountably many values even though a parameter-free finite syntax has only countably many strings.

For computable selectors, a program for $A$ provides a finite effective input to the normal-form construction. For noncomputable but finitely definable selectors, one obtains a finite semantic description without an effective copy. For selectors outside a chosen finite-description language, the family formula still proves existence relative to that selector, but does not internalize it into that language.

These are exact distinctions between a construction schema, its parameters, and an individual realized object. They are the appropriate replacement for an absolute label such as ``non-Komorebi group.''

\chapter{A worked arithmetic atlas}
\label{ch:atlas}
\chapterintro{The examples below retain the infinite definition while calculating only the finite coordinates that a particular operation needs. They make the cocycle and the root certificates visible.}

\section{The empty selector}
Take $A=\varnothing$. The selected primes are
\[
S_A=\{2,5,11,17,23,31,41,47,\ldots\},
\]
namely the even-indexed primes. The first coordinates of three rational lifts are
\begin{center}
\begin{tabular}{@{}crrrrrr@{}}
\toprule
$p$&2&3&5&7&11&13\\\midrule
$u_A=s_A(1)$&1&0&1&0&1&0\\
$s_A(1/2)$&0&0&3&0&6&0\\
$s_A(1/3)$&1&0&2&0&4&0\\\bottomrule
\end{tabular}
\end{center}
Each residue is exact. At $p=2$, the lift of $1/2$ is set to zero because the denominator is not invertible. At $p=5$, its value is $3$ because $2\cdot3=1$ in $\F_5$.

Doubling the half-lift agrees with $u_A$ at every prime except $2$. Therefore
\[
2(1/2,0)=(1,e_2).
\]
The element $u_A$ has no half in the group: there is no $y$ with $2y=u_A$, because its $2$-coordinate is nonzero. The apparent candidate $(1/2,0)$ misses by precisely the certificate coordinate.

At the unselected prime $3$, the coordinate of $u_A$ is zero. In this case
\[
3(1/3,0)=(1,0).
\]
The three roots are
\[
(1/3,0),\qquad(1/3,e_3),\qquad(1/3,2e_3).
\]
Their differences form $G_A[3]=C_3$, as predicted by Proposition~\ref{thm:rootfiber}.

\section{A successful composite-root calculation}
Let
\[
g=s_A(1)+e_2=(1,e_2).
\]
Its coordinate at $2$ is $1+1=0$, and its coordinate at $3$ is zero. Thus the multiplier $6$ has no obstructing coordinate.

The rational candidate is $t=1/6$. At a selected prime other than $2$, multiplication by $6$ recovers the coordinate of $u_A$ because the denominator is invertible, except possibly at $3$, which is unselected. At $2$, both sides are zero. Hence
\[
6(1/6,0)=(1,e_2)=g.
\]
Every other root is obtained by adding an arbitrary element of $C_2\oplus C_3$. There are exactly six roots. This calculation exhibits a positive certificate without searching an unbounded domain.

Now contrast $u_A$ itself. The same multiplier $6$ fails at the prime $2$. No change at the $3$-coordinate and no choice of rational representative can repair that obstruction while keeping the target fixed.

\section{A torsion correction that cannot be discarded}
Take the purely torsion target $e_3$ and multiplier $2$. The root criterion succeeds because the target has zero $2$-coordinate. Its rational part is zero, but the answer is not the zero element:
\[
2(0,2e_3)=(0,e_3).
\]
In $\F_3$, multiplication by $2$ sends $2$ to $1$. This is the simplest example showing why the finite correction in the root algorithm is necessary even when the rational coordinate causes no difficulty.

\section{Finite quotient representatives}
For any target $g$ and positive multiplier $m$, define the finite vector
\[
d_m(g)=\sum_{p\mid m}g_p e_p.
\]
Then $g-d_m(g)$ has zero coordinates at every prime dividing $m$. The root theorem provides a witness that it belongs to $mG_A$. Thus $d_m(g)$ is an exact finite representative of the coset of $g$.

For $m=12$, only coordinates $2$ and $3$ matter. The six representatives are
\[
0,\ e_2,\ e_3,\ 2e_3,\ e_2+e_3,\ e_2+2e_3.
\]
The vector $e_2+e_3$ generates this quotient, which is $C_6$, not $C_{12}$. Repeated prime powers in the multiplier do not create additional coordinates.

\begin{figure}[tbp]
\centering\includegraphics[width=\textwidth]{\figfile{07_squarefree_quotients.pdf}}
\caption{The exact order of $G_A/mG_A$ for $1\le m\le30$ is $\rad(m)$. The values are independent of the selector. This plot evaluates the theorem's arithmetic formula, rather than estimating the size of an unknown quotient.}
\label{fig:radical}
\end{figure}

\section{Changing one selector bit}
Let $B=\{0\}$, so the first pair changes from selecting $2$ to selecting $3$. The groups $G_A$ and $G_B$ are equal as subgroups of $P$. Their half-lift coordinates differ, however:
\[
s_A(1/2)-s_B(1/2)=e_3.
\]
At $p=2$, both half-lifts are zero. At $p=3$, the $B$ lift is $2$, so its negative is $1$. Therefore the same ambient element has normal forms
\[
(1/2,0)_A=(1/2,e_3)_B.
\]
Doubling the right-hand side in the $B$ group gives
\[
(1,2e_3)_B,
\]
which is the transported form of $(1,e_2)_A$. The correction in \eqref{eq:transport} is required to compare the two group laws. Merely reusing the same pair $(r,d)$ under a new selector changes the represented element.

\section{A finite tail exception}
Consider an infinite-order element
\[
x=s_A(5/6)+2e_7.
\]
An exception set for its signature is $\{2,3,5,7\}$. Outside these primes, the denominator is invertible, the numerator is nonzero, and the torsion correction vanishes. Thus the root pattern of $x$ agrees with that of $u_A$ in every later prime pair.

This example illustrates the nonuniform step in Theorem~\ref{thm:degree}. In known normal-form coordinates the exception set is computable. In an arbitrary unmarked copy, a reduction may use a finite exception bound without uniformly discovering it from the copy.

\section{Reproducing the examples}
From the package root, the following uses the supplied standard-library implementation:
\begin{lstlisting}[language=Python]
from fractions import Fraction
from src.dense_lift_groups import DenseLiftGroup, Element

G = DenseLiftGroup(lambda n: False)
u = Element.make(1)
half = Element.make(Fraction(1, 2))

assert G.scale(2, half) == Element.make(1, {2: 1})
assert G.root(u, 2).obstructing_prime == 2
third = G.root(u, 3).root
assert third is not None and G.scale(3, third) == u

x = Element.make(1, {2: 1})
sixth = G.root(x, 6).root
assert sixth is not None and G.scale(6, sixth) == x
\end{lstlisting}
Only computable selector examples are executed. A callable supplied as an oracle is assumed to return the correct selector bit. The implementation does not pretend to evaluate the halting set or an arbitrary noncomputable parameter.

\chapter{Boundaries as part of the calculus}
\label{ch:calculus}
\chapterintro{The constructions identify specific obligations for realization: compatibility, admissibility, generation, and effective evaluation. A failure at one level need not be a failure at another.}

\section{The ambient object and its admissible realizations}
The chain
\[
D\subseteq G_A\subseteq P
\]
contains three different mathematical roles. The direct sum provides the finite-support corrections. The product supplies every compatible prime-coordinate tuple. The intermediate group is the intended admissible realization, determined by a rational line in the quotient.

This arrangement gives a concrete model for a central question of Komorebi Space. Once the compatible ambient object exists, what additional data identifies the realization under study? Here the answer is exact: a rational subspace of $P/D$, or in the paired family its selector class. Existence of the ambient product does not eliminate that choice.

The construction does not assert that every Komorebi Space model is a prime-coordinate product. It isolates a theorem-level phenomenon that any broad realization language must be able to express: the space of compatible observations can be larger than the intended source image.

\section{Three certificates and one global obstruction}
The finite-seed argument supplies a projection certificate. An omitted prime coordinate vanishes on every admissible endomorphism orbit but not on the target group. This is a structural failure of generation under the fixed operator action.

The root algorithm supplies a coordinate certificate of impossibility or a finite correction witnessing success. It is a fully local algebraic realization problem. The successful and unsuccessful branches are both explicit.

The nonsplitting result supplies an extension obstruction. A rational quotient exists, but no homomorphic section realizes it inside the reduced product. In cocycle coordinates, the obstruction is the impossibility of removing the defect by a $D$-valued correction.

The grammar theorem is different. It excludes every finite description in a fixed language by an infinite diagonal comparison. It need not supply a terminating rejection procedure for individual candidate descriptions. Conflating this global semantic obstruction with the finite root certificates would overstate its algorithmic content.

\section{A realization profile rather than a single score}
One scalar cannot capture all the boundaries exhibited here. The following profile retains the distinctions:
\[
\bigl(\text{seed rank},\ \text{grammar complexity},\
\text{least presentation degree},\ \text{observation type}\bigr).
\]
For every $G_A$, the seed rank is $\aleph_0$ and the observation type is the same. The least presentation degree varies with $A$. Grammar complexity depends on which selectors the chosen semantic language can name.

The profile is not proposed as a complete classification invariant for arbitrary mathematical objects. It is a disciplined way to avoid substituting one question for another. A theory may have exact coordinates without a finite generator, a finite formula without computable equality, or an effectively presented inverse-limit completion without a computable copy of the intended dense realization.

\section{What a machine-facing contract should retain}
A reusable construction record should specify its input type, its permitted parameters, its output object, and the sense in which the output is determined. It should also state whether it returns a concrete representation, an abstract isomorphism type, or an existence proof. When completion is used, the intended image and the completed ambient object should be separate fields.

For the present family, a compact exact record is:
\begin{center}
\begin{tabularx}{\textwidth}{@{}p{0.22\textwidth}X@{}}
\toprule
Field&Content\\\midrule
Ambient system&$D\subseteq P$ and the quotient $q:P\to P/D$.\\
Parameter&The selector $A$, with its declared oracle or finite semantic description.\\
Admissible image&$q^{-1}(\Q q(u_A))$.\\
Concrete names&Reduced rationals paired with normalized finite-support vectors.\\
Operation&The finite cocycle law \eqref{eq:grouplaw}.\\
Positive certificate&A root and its exact multiplication check.\\
Negative certificate&A prime coordinate obstructing the required root.\\
Global invariants&Selector modulo finite changes and least presentation degree.\\
Observation boundary&Unmarked finite quotients forget the admissible source image.\\\bottomrule
\end{tabularx}
\end{center}
Such a record does not solve arbitrary group-theoretic problems. It makes the semantics of this construction inspectable and prevents a machine from promoting a finite observation theorem into a source reconstruction theorem.

\section{Further questions with exact starting points}
The general correspondence of Chapter~\ref{ch:purity} extends the ambient construction to higher-dimensional rational subspaces $W\le P/D$. Its finite quotient theorem is already complete. What remains is to classify those subspaces under the diagonal multipliers that induce group homomorphisms. The paired rank-one family has an exact answer, but the book does not claim that its binary signature classifies higher rank.

A second question concerns effective selection. For $G_A$, the normal-form operation and the paired root tests give a matching upper and lower bound. For a more general countable $W$, one must specify how its rational dependencies are presented and how lifts are computed. A list of candidate spanning vectors is not enough if their dependence relations are unavailable.

A third question concerns enriched observations. A coherent marked point restores the selector in the declared rank-one family. A useful general theorem would characterize which additional marks or admissibility predicates suffice for reconstruction in larger classes. Such a result must state whether its reconstruction is semantic, effective, or functorial.

Finally, the formal companion keeps the arithmetic kernel, quotient lemmas, tail signature, and computability reduction as distinct proof obligations. Appendix~\ref{app:formal} explains the resulting interfaces. The oracle theorem has its own effective semantics, with native oracle programs, computable equality, and terminating searches in addition to the algebraic proofs.

\section{Conclusion}
The principal boundary is now concrete. There is a fixed ambient completion, presented by a computable inverse system, and a family of countable, residually finite groups densely realized inside it. Their finite-quotient systems are canonically identified, and restriction to $D$ naturally identifies their homomorphism sets into every finite target. Yet the groups have distinct algebraic types and require arbitrarily different computational information.

Realization is therefore not exhausted by compatibility. It includes the selection of admissible objects, the transport of their meaning, and an account of the information needed to evaluate them. The groups in this book turn that distinction into explicit constructions, exact classifications, and reproducible finite certificates.

\appendix
\chapter{Exact computation and reproduction}
\label{app:implementation}
\newcommand{\BaseAssertions}{114,738}
\newcommand{\ExtraAssertions}{53,284}
\newcommand{\ReleaseAssertions}{21,047}
\newcommand{\AllAssertions}{189,069}

\chapterintro{The computational companion evaluates the finite arithmetic of the construction. It does not stand in for the infinite proofs. This appendix states its contracts, its checks, and its limits.}

\section{A finite representation of each named element}
An element is stored as a reduced rational and a normalized finite-support tuple of prime residues. The tuple is ordered by prime label, has no repeated labels, and omits zero residues. Use the constructor \texttt{Element.make} to obtain this canonical form. Direct construction enforces the same normalization. Rational inputs must be integers or \texttt{Fraction} values. Residues, prime labels, and multipliers must be integers. Floating-point values are rejected rather than rounded or silently truncated.

The selector is supplied as a callable taking a nonnegative integer $n$ to the Boolean value of $\one_A(n)$. Integer bits zero and one are also accepted. Other return values are rejected. The group object converts a prime label to its index, determines which pair it belongs to, and queries the corresponding selector bit. This is an oracle interface. A finite Python function implements only those selectors for which it can return the correct value. No program in the package claims to compute a noncomputable selector.

\begin{center}
\begin{tabularx}{\textwidth}{@{}p{0.35\textwidth}X@{}}
\toprule
Interface&Contract\\\midrule
\texttt{Element.make(r, d)}&Normalize a rational and finite torsion data.\\
\texttt{section\_value(r, p)}&Evaluate one coordinate of the specified section.\\
\texttt{cocycle(r, t)}&Return the exact finite correction for rational addition.\\
\texttt{add(x, y)}&Apply the cocycle group law.\\
\texttt{scale(m, x)}&Integer multiplication by repeated doubling.\\
\texttt{root(x, m)}&For $m\ge1$, return a root or an obstructing prime.\\
\texttt{finite\_change\_transport}&Change normal forms when a complete finite difference set is supplied.\\\bottomrule
\end{tabularx}
\end{center}

The finite-change transport routine has a global precondition: the selectors differ only at the supplied indices. No finite examination of arbitrary oracle callables can verify that precondition. The executable tests use selectors constructed to satisfy it exactly.

\section{Root construction as an algorithm}
The proof of Theorem~\ref{thm:rootcertificate} gives the following procedure.
\begin{enumerate}
\item Factor the positive integer $m$ and inspect the target coordinates at its prime divisors. If any is nonzero, return that prime as the obstruction.
\item Divide the rational coordinate by $m$ to obtain a candidate section lift.
\item Compute the finite residual between the target and $m$ times that lift. Its support is confined to the input torsion support and the denominator primes stated in the proof.
\item At primes not dividing $m$, divide the residual coordinate by $m$ in the finite field. At primes dividing $m$, the residual is zero. Add this finite correction to the candidate lift.
\end{enumerate}
The procedure returns a complete normal-form element. Its certificate is checked by the independently implemented integer multiplication routine. The negative certificate is checked by verifying both $p\mid m$ and a nonzero target coordinate.

Nonpositive or noninteger root multipliers are rejected. The routine does not enumerate all elements of the group. Its finite support bound comes from the theorem. This distinguishes a certified construction from an unbounded search that might fail to terminate when no root exists.

\section{What the tests cover}
The core suite checks cocycle symmetry, support, and coordinate formulas over a specified finite rational set. It checks the cocycle identity on all triples from that set for eight periodic selector masks. It also tests group identities on reproducibly generated exact elements, root witnesses and obstructions, quotient representatives, marked bit decoding, tail signatures outside explicit exceptions, and finite-change transport.

The supplemental suite checks the monograph's additional calculations. It enumerates entire finite root fibers and verifies their cardinalities. It checks explicit diagonal endomorphisms against coordinate multiplication and addition, and reproduces the worked sixth-root, torsion-root, and coordinate-change examples.

The verification run passed \BaseAssertions{} core assertions, \ExtraAssertions{} supplemental assertions, and \ReleaseAssertions{} release-regression assertions, for a total of \AllAssertions{} exact finite assertions. Seven named mathematical negative controls reject the false simplifications below. The release suite additionally checks malformed-input rejection, normal-form normalization, and diagonal automorphism inversion. Run reports record the executed source hashes and parameters.

\begin{center}
\begin{tabularx}{\textwidth}{@{}p{0.39\textwidth}X@{}}
\toprule
False simplification rejected&Reason for rejection\\\midrule
Omit the cocycle&Twice the half-lift can differ from $u_A$ at prime $2$.\\
Every rational division lifts&A nonzero divisor-prime coordinate obstructs a root.\\
Reuse coordinates after a selector change&The same pair may represent a different ambient element.\\
Discard torsion correction in a root&A pure torsion target may need a nonzero finite correction.\\
Use $C_m$ for every integer quotient&Prime powers collapse to $C_{\rad(m)}$.\\
Allow a zero root multiplier&The root contract requires a positive integer.\\
Treat rational diagonal representatives as a ring embedding&Exceptional denominator coordinates can destroy multiplicativity of a chosen section.\\\bottomrule
\end{tabularx}
\end{center}

These checks exercise intended failure mechanisms. They are not a random test of historical originality or an attempt to compute a halting oracle. The universal results about all selectors, isomorphism classes, and Turing degrees are supported by the printed proofs.

\section{Reproduction commands}
Python 3.10 or later is sufficient for the exact arithmetic. From the package root, run
\begin{lstlisting}
python src/dense_lift_groups.py
python src/exact_checks.py --output verification/local_core.json
python src/supplemental_checks.py --output verification/local_extra.json
python src/release_checks.py --output verification/local_release.json
\end{lstlisting}
The tests use only the standard library. Matplotlib is needed only to regenerate the vector figures:
\begin{lstlisting}
python src/make_figures.py
\end{lstlisting}
Rebuild the monograph with \texttt{python build.py}. The README lists the required TeX packages. Once dependencies are installed, the build uses no network access. The package retains the editable manuscript, figures, tests, theorem map, and original research archives.

\section{Verification boundary}
The separate Lean companion verifies the numbered mathematical contracts, including the infinite classification, oracle criteria, and grammar escape. Appendix~\ref{app:formal} explains its interfaces and trust boundary. Finite Python tests and typesetting checks support the implementation and publication. They are recorded separately from the compiler, declaration audit, and kernel replay receipts.

\chapter{Formal companion and proof guide}
\label{app:formal}
\chapterintro{The formal companion works with the actual infinite groups and their canonical completions. Its computational theorems use explicit oracle programs and decidable equality on named copies. This appendix explains the interfaces and the evidence behind them.}

\section{The proof dependencies}
\begin{figure}[htbp]
\centering\includegraphics[width=\textwidth]{\figfile{08_proof_dependencies.pdf}}
\caption{Principal dependencies of the proof development. Finite arithmetic checks illustrate the constructions. The infinite theorems are proved independently, and effective conclusions use their own oracle and presentation semantics.}
\end{figure}

The companion retains both the ambient preimage and the finite normal form. The ambient model supplies the purity, homomorphism, and completion arguments. The finite model supplies executable arithmetic. Coordinate identities connect the two models, so a finite calculation is never substituted for a theorem about the infinite object.

The machine readable statement registry records all numbered contracts in the book, their Lean declarations, and their exact scope. The table below lists the principal dependencies. Compiler receipts and source hashes accompany the companion rather than interrupting the mathematical exposition.

\section{Statement registry}
\small
\begin{longtable}{@{}>{\raggedright\arraybackslash}p{0.16\textwidth}>{\raggedright\arraybackslash}p{0.45\textwidth}>{\raggedright\arraybackslash}p{0.31\textwidth}@{}}
\toprule
Statement&Conclusion&Key dependency\\\midrule
\endfirsthead
\toprule Statement&Conclusion&Key dependency\\\midrule\endhead
\ref{thm:reach}&An inaccessible packet obstructs generation.&Typed block-support condition.\\
\ref{thm:torP}&$\Tor(P)=D$, and $D$ is dense.&Distinct prime coordinates.\\
\ref{thm:endD}&All endomorphisms of $D$ are diagonal.&Prime annihilators.\\
\ref{thm:rankD}&$D$ has infinite endomorphism-seed rank.&Finite support.\\
\ref{thm:procyclic}&$P$ is topologically cyclic.&Chinese remainder theorem.\\
\ref{thm:Dcompletion}&$\widehat D=P$.&Cofinal finite coordinate kernels.\\
\ref{thm:Vrational}&$P/D$ is a rational vector space.&Unique integer division modulo $D$.\\
\ref{thm:subD}&$\Sub(D)\cong\mathcal P(\N)$.&Coordinate isolation by CRT.\\
\ref{thm:spectrumce}&Computable copies of $D_B$ correspond to enumerable selectors.&Prime-order witness searches.\\
\ref{thm:unary}&Unary selector groups correspond to ultimately periodic sets.&Finite-state orbit.\\
\ref{thm:purecorrespondence}&Pure intermediate groups are rational preimages.&Finite adjustment at divisor primes.\\
\ref{thm:generalquotients}&All preimages $H_W$ have completion $P$.&Purity and finite-target exponent.\\
\ref{thm:finitehom}&Restriction preserves all maps to finite targets.&Finite coordinate factorization.\\
\ref{thm:normalform}&$G_A$ has unique $(r,d)$ names.&Nonzero rational line.\\
\ref{thm:cocyclesupport}&Each cocycle value has controlled finite support.&Reduction outside denominator primes.\\
\ref{thm:cocyclelaw}&The normal forms form the actual group $G_A$.&Cocycle identity.\\
\ref{thm:cocyclebit}&A cocycle value reads one selector coordinate.&The chosen section at one prime.\\
\ref{thm:rootcertificate}&A root or a prime obstruction is constructible.&Finite residual correction.\\
\ref{thm:rootfiber}&A nonempty root fiber has $\rad(m)$ elements.&The prime-coordinate kernel.\\
\ref{thm:quotientformula}&$G_A/mG_A=C_{\rad(m)}$.&Root kernel and finite surjectivity.\\
\ref{thm:tailsignature}&Every nontorsion element has the same tail signature.&Finite numerator, denominator, and support exceptions.\\
\ref{thm:isoclass}&$G_A\cong G_B$ iff $A\eqfin B$.&Intrinsic divisibility and finite transport.\\
\ref{thm:allhom}&Every homomorphism is a diagonal multiplier.&The quotients $G_A/pG_A$.\\
\ref{thm:homobstruction}&Different eventual classes admit only torsion-image maps.&One selected prime per pair.\\
\ref{thm:endGA}&Exact endomorphism-ring description.&Diagonal action on the rational line.\\
\ref{thm:rankGA}&Every $G_A$ has infinite seed rank.&Infinitely many omitted unselected coordinates.\\
\ref{thm:degree}&An oracle computes a copy iff it computes $A$.&Cocycle evaluation and complementary root searches.\\
\ref{thm:nodecoder}&No recovery from the unmarked quotient system.&Two nonisomorphic common-completion groups.\\
\ref{thm:faithful}&Completion is faithful but not full on the paired family.&Finite separation and the torsion-image obstruction.\\
\ref{thm:rankonereflection}&Completion reflects isomorphisms between paired-prime groups.&Diagonal units and the exact source category.\\
\ref{thm:escape}&Continuum many groups escape a fixed countable grammar.&Repeated diagonal disagreement.\\
\bottomrule
\end{longtable}
\normalsize


\section{Reading the formal layers}
\subsection{Ambient groups, subgroups, and generation}
The prime product is an unrestricted dependent function space. Its finite support subgroup is defined inside that product. The quotient becomes a rational vector space only after positive integer multiplication is proved bijective on the actual quotient.

The subgroup classification isolates coordinates by Chinese remainder interpolation and reconstructs every finite support vector from prime seeds. Intrinsic prime torsion gives abstract isomorphism classification. The topology on the subgroup lattice is independently induced by membership functions, and the resulting Cantor homeomorphism is then proved continuous in both directions.

The general reachability theorem uses a genuine dependent direct sum, primitive transfer blocks, finite operator words, and ordinary finite directed paths. An omitted nonzero packet supplies both the algebraic obstruction and, under the stated continuity hypotheses, the closed kernel obstruction to topological generation.

\subsection{Finite arithmetic and roots}
The executable root solver takes a rational code, a finite residue list, and a selector oracle. Each list position names a prime value. Positions that are not prime do not contribute to the interpreted vector. Consequently raw list equality is not the equality of represented group elements.

Every finite support correction admits a residue list, and every element of the original selector group admits one of the finite normal form codes. A successful solver output is proved to satisfy the root equation at every prime. A rejected request returns a prime dividing the multiplier with a nonzero target coordinate. The correction is zero at divisor primes. This normalization concerns the correction, and need not make the entire root coordinate zero.

The traced implementation erases to the executable solver. If $B$ is its finite correction bound, it uses at most $m+1+2B$ selector queries and at most $m+1+B$ loop candidates. These are explicit oracle and iteration counts. They are not bounds on the bit cost of rational arithmetic, primality testing, division, or modular inversion. The Lean proofs certify the Lean implementation and its mathematical interpretation. They do not claim kernel verification of the separate Python runtime.

\subsection{Computable copies and degrees}
The effective layer uses native oracle recursion with oracle calls, primitive recursion, and terminating searches. A named infinite copy has a bijection from the natural numbers, a specified zero, and an oracle computable addition operation. Equality of names is decidable. Inversion and subtraction are computed by proved terminating searches.

Redundant finite codes are permitted only after equality of their represented elements is proved oracle computable. The companion then constructs least representatives, computes their enumeration, and obtains genuinely bijective natural names. The concrete normal form supplies all those hypotheses through rational arithmetic and finite cocycle evaluation.

Marked decoding performs two complementary positive root searches. For an unmarked copy, the existence proof may choose a finite natural name of the distinguished element under an abstract isomorphism and hard code that name. It does not uniformly locate the distinguished element from an arbitrary copy. This gives the same existential oracle criterion while preserving the marked and unmarked distinction.

The degree statements use the native Turing reducibility relation. Ordinary degrees of sets are represented by total binary membership oracles. The least degree is attained by a constructed copy and is a lower bound for every other oracle computing a copy. All members retain the same finite quotient observation system.

\subsection{Observations, categories, and fixed grammars}
The observation interface contains finite groups, transition maps, and their cofiltered index category. It does not contain the source group or evaluations of named source points. The common coordinate subsystem is proved categorically cofinal in the system of all finite quotient kernels. A natural isomorphism identifies its reindexed diagrams. The no decoder theorem therefore concerns the actual unmarked quotient systems.

Faithfulness uses the official completion functor on all residually finite groups. Failure of fullness is packaged on the paired family. Failure to reflect isomorphisms is packaged on the larger full category that also contains the direct sum. These are distinct source categories, as required by the mathematical statements.

Grammar escape uses infinitely repeated disagreement against every semantic output. Independent repeated bits retain continuum many pairwise nonisomorphic groups, with no embeddings between distinct members. The semantics may be noncomputable. The theorem excludes outputs of the fixed grammar and does not assert escape from every possible surrounding language.

\section{Reproduction and trust}
The companion includes source files, exact statement contracts, positive examples, rejection controls, source hashes, and compiler receipts. The verification script rebuilds the production modules with warnings treated as errors, audits their declaration dependencies, executes the contracts, and replays the project modules with the supplied kernel checker. The portable runtime and dependency revisions are pinned in the accompanying machine readable files.

The rejection controls test specific false strengthenings, including universal root existence, the wrong cyclic quotient modulus, literal unmarked selector recovery, removal of required finiteness hypotheses, and reconstruction from the unmarked observation system. Each rejection must contain the intended diagnostic. A missing import or an unrelated elaboration error is not accepted as a mathematical rejection.

Only the standard logical dependencies recorded by the companion are admitted. Imported dependency caches are trusted inputs whose bytes are checked during the run. Kernel replay uses the same Lean implementation, so it is not an independent kernel or external peer review. The source archive and its receipts preserve enough information to repeat the checks within that stated trust boundary.

\chapter{Glossary}
\label{app:glossary}
\markboth{Glossary}{}
\begin{description}[style=nextline,leftmargin=0pt,labelindent=0pt,labelwidth=0pt,labelsep=0pt,itemindent=0pt,font=\bfseries\color{Ink},itemsep=9pt]
\Needspace{5\baselineskip}\item[Admissible realization] The particular object or image allowed by a construction contract, as distinct from every compatible tuple in an ambient completion.
\Needspace{5\baselineskip}\item[Atomic diagram] The effective record of the basic relations and operations of a countable named structure. It is the standard object used to measure the degree of a copy.
\Needspace{5\baselineskip}\item[Cocycle] Here, a symmetric map $c:\Q\times\Q\to D$ satisfying the identity that makes the twisted addition on $\Q\times D$ associative.
\Needspace{5\baselineskip}\item[Cofinal quotient system] A family of finite quotients through which every finite quotient factors. Its inverse limit computes the full profinite completion.
\Needspace{5\baselineskip}\item[Computable copy] A countable realization with effective names, decidable equality, and computable operations. An oracle version permits the stated oracle.
\Needspace{5\baselineskip}\item[Computable profinite presentation] An effective inverse system of finite groups and surjective transition maps whose limit is the given profinite group. This is not a computable enumeration of all elements of an uncountable product.
\Needspace{5\baselineskip}\item[Computably enumerable set] A set whose members can be listed by a procedure. Nonmembership need not be decidable.
\Needspace{5\baselineskip}\item[Completion] An ambient object obtained by admitting compatible limits in a specified topology. In this book the relevant completion is profinite.
\Needspace{5\baselineskip}\item[Dense subgroup] A subgroup meeting every nonempty basic open set. Density does not imply equality with the ambient group.
\Needspace{5\baselineskip}\item[Divisible abelian group] An abelian group in which every equation $my=x$ with $m\ge1$ is solvable. A divisible quotient need not split off as a subgroup.
\Needspace{5\baselineskip}\item[Endomorphism-seed rank] The minimum number of elements generating a group as a module over its full endomorphism ring.
\Needspace{5\baselineskip}\item[Eventual equality] Equality modulo finitely many coordinates, written $A\eqfin B$. This is the unmarked selector invariant of $G_A$.
\Needspace{5\baselineskip}\item[Finite description] A valid finite code in a specified language with declared semantics. It need not provide computable operations.
\Needspace{5\baselineskip}\item[Finite support] A vector has only finitely many nonzero coordinates. Such vectors form $D$ inside $P$.
\Needspace{5\baselineskip}\item[Fully invariant subgroup] A subgroup preserved by every endomorphism of its ambient group. Every subgroup of the squarefree direct sum has this property.
\Needspace{5\baselineskip}\item[Grammar-relative nonrealizability] Failure to have an admissible description in one fixed grammar. It is not absolute mathematical nonexistence or indescribability.
\Needspace{5\baselineskip}\item[Marked group] A group together with named elements that isomorphisms must preserve. The mark $u_A$ retains the literal selector.
\Needspace{5\baselineskip}\item[Normal form] The unique pair $(r,d)$ representing $s_A(r)+d$. Equality of these names is simpler than the selector-dependent addition law.
\Needspace{5\baselineskip}\item[Oracle] An external set-membership resource available to a computation. Its information is charged to the effective realization contract.
\Needspace{5\baselineskip}\item[Prime spectrum] The prime orders of nonzero torsion elements. It classifies the groups $D_B$ but is constant across all $G_A$.
\Needspace{5\baselineskip}\item[Profinite completion] The inverse limit of all finite quotient groups, denoted $\widehat G$. It is not the Pontryagin dual.
\Needspace{5\baselineskip}\item[Procyclic group] A topological group with a dense cyclic subgroup. The product $P$ is procyclic.
\Needspace{5\baselineskip}\item[Pure subgroup] A subgroup $H\le M$ for which $H\cap mM=mH$ for all positive integers $m$.
\Needspace{5\baselineskip}\item[Rational line] A one-dimensional subspace over $\Q$. The preimage of $\Q q(u_A)$ is $G_A$.
\Needspace{5\baselineskip}\item[Reduced group] A group with no nonzero divisible subgroup. This property does not exclude divisible quotients.
\Needspace{5\baselineskip}\item[Residual finiteness] Every nonzero element is detected by some finite quotient. It does not imply source reconstruction from the unmarked quotient system.
\Needspace{5\baselineskip}\item[Root certificate] Either an exact element $y$ with $my=g$, or a prime divisor of $m$ at which the target has a nonzero coordinate.
\Needspace{5\baselineskip}\item[Selector] A subset $A\subseteq\N$ whose bits choose one prime from each consecutive pair. It determines the admissible rational line.
\Needspace{5\baselineskip}\item[Squarefree radical] The integer $\rad(m)$ formed by retaining each prime divisor of $m$ once.
\Needspace{5\baselineskip}\item[Tail signature] The pattern of prime-root solvability outside a finite exception set. Every infinite-order element of $G_A$ has the selector's tail signature.
\Needspace{5\baselineskip}\item[Torsion] Elements of finite order. In the product and in every $G_A$, the torsion subgroup is exactly $D$.
\Needspace{5\baselineskip}\item[Turing degree] A measure of relative computability under mutual oracle computation. The least presentation degree of $G_A$ is the degree of $A$.
\Needspace{5\baselineskip}\item[Unmarked observations] Finite quotient objects and transition maps without a named coherent source element or the original admissible image.
\end{description}

\backmatter
\begin{thebibliography}{9}
\addcontentsline{toc}{chapter}{References}
\markboth{References}{}
\bibitem{Montalban}
Antonio Montalb\'an.
\emph{Computable Structure Theory: Within the Arithmetic}.
Perspectives in Logic, Cambridge University Press, 2021.
For the cited numbering, see the author's draft of June 25, 2019, Example III.39 and the following paragraph, p.~50.
\url{https://math.berkeley.edu/~antonio/CSTpart1.pdf}.

\bibitem{Calvert}
Wesley Calvert, Valentina Harizanov, and Alexandra Shlapentokh.
Turing degrees of isomorphism types of algebraic objects.
\emph{Journal of the London Mathematical Society} 75(2) (2007), 273--286.
\href{https://doi.org/10.1112/jlms/jdl012}{doi:10.1112/jlms/jdl012}.
Preprint: \href{https://arxiv.org/abs/math/0507128}{arXiv:math/0507128}.

\bibitem{Fuchs}
L\'aszl\'o Fuchs.
\emph{Abelian Groups}.
Springer Monographs in Mathematics, Springer, 2015.
\href{https://doi.org/10.1007/978-3-319-19422-6}{doi:10.1007/978-3-319-19422-6}.
General background for the terminology of pure subgroups and mixed abelian groups. The arguments used here are proved in the text.

\bibitem{Lupini}
Martino Lupini, Alexander Melnikov, and Andr\'e Nies.
\emph{Computable topological abelian groups}.
\href{https://arxiv.org/abs/2105.12897v2}{arXiv:2105.12897v2}, 2021.
Definition 3.3 supplies the computable inverse-system convention. Section 3.2 describes symmetric cocycles and abelian extensions.

\bibitem{Wilkes}
Gareth Wilkes.
\emph{Part III Profinite Groups}.
Lecture notes, Lent Term 2020.
Sections 1.2 and 3.1 discuss profinite completion and its functoriality.
\url{https://www.dpmms.cam.ac.uk/~grw46/LectureNotes.pdf}.
\end{thebibliography}

\end{document}
